% [原文 p.365]
\chapter{S 类理论中的 BPS 态}

\ChapterAbstract{我们综述四维 $\mathcal{N}=2$ 理论中 BPS 态的几何方法，重点是 S 类理论。由于这部分材料较为高深，我们的处理方式是对文献作导引式综述，而非全面导论。我们先按 [151, 155] 把 BPS 态描述为在 S 类理论相伴的紫外曲线 $C_{\mathrm{UV}}$ 上具有变张力的有限弦网，并引入 $A_1$ 型理论情形下的 WKB 三角剖分概念。然后转向谱网络框架——它由 [154] 引入，建立在 [153, 152] 发展的 $2d$–$4d$ 带标架 BPS 态理论基础之上。}

\section{变张力弦}

\subsection{来自 M 理论的弦}

在接下来的讨论中，我们将特别注意 $A_{N-1}$ 型的 S 类理论，它们由 $A_{N-1}$ 型相互作用 $6d$ $\mathcal{N}=(2,0)$ SCFT 在带穿孔的黎曼曲面上紧化而构造。我们会经常借助 M 理论的视角：在那里 $6d$ 理论作为平坦空间 $\R^{1,10}$ 中一叠 $N$ 个 M5-膜上的世界体理论出现。我们的叙述遵循 [155]。

与 D-膜的情形类似，这些 $6d$ $\mathcal{N}=(2,0)$ SCFT 表现出一种库仑分支，在其上动力学一般约化为秩为 $N-1$ 的阿贝尔 $6d$ $\mathcal{N}=(2,0)$ SCFT。在 M 理论中，这个库仑分支对应于把各 M5-膜在其世界体的五个横向方向上分开。这一分离是通过给 $A_{N-1}$ 型 $6d$ 理论中特定的手征算子赋真空期望值（VEV）来实现的。

在其库仑分支上，$A_{N-1}$ 型的 $6d$ $\mathcal{N}=(2,0)$ SCFT 支持由整数对 $(ij)$（其中 $1\le i\ne j\le N$）标号的 BPS 弦。在 M 理论图像中，这些 BPS 弦表现为在第 $i$ 个与

% [原文 p.366]
第 $j$ 个 M5-膜之间张开的 M2-膜（图 22.1）。一条 $(ij)$ BPS 弦的张力正比于第 $i$ 个与第 $j$ 个 M5-膜之间的间距。

\Fig[0.95\linewidth]{figures/fig_22_1.png}{22.1}{$A_2$ 型 $6d$ $\mathcal{N}=(2,0)$ SCFT 中的弦。左图为 M 理论图像，显示在 M5-膜之间张开的 M2-膜。右图显示 $6d$ $\mathcal{N}=(2,0)$ 理论中相应的弦。}

$(ij)$、$(jk)$ 与 $(ik)$ 型的弦可以在接点处相遇，只要满足保持超对称的特定条件：相应的 M2-膜接点必须满足「力学平衡」。

现在转向 $A_{N-1}$ 型的 S 类理论。在上一节中我们看到，这些理论的库仑分支在 M 理论中被描述为一叠 $N$ 个环绕紫外曲线 $C_{g,n}$ 的 M5-膜的可归一化形变空间。更确切地说，所考虑的设定如下：我们考虑 $\R^{1,6}\times Q$ 上的 M 理论，其中 $Q$ 是超凯勒流形，$N$ 个 M5-膜环绕 $Q$ 中的复曲线 $C_{g,n}\subset Q$。理论 $\mathcal{S}[A_{N-1},C_{g,n},\mathcal{D}]$ 的库仑分支对应于 M5-膜叠在 $Q$ 中的可归一化连通形变 $\Sigma$，它们构成 $C_{g,n}$ 的 $N$ 次分歧覆叠。缺陷数据 $\mathcal{D}$ 规定了 $\Sigma$ 在 $C_{g,n}$ 的穿孔处的边界条件。

$C_{g,n}$ 在 $Q$ 中的邻域可以等同于 $C_{g,n}$ 的全纯余切丛 $T^*_{1,0}(C_{g,n})$。$C_{g,n}$ 附近的激发对应于 M 理论构型的场论极限。因此，S 类理论的库仑分支可以等同于 $T^*_{1,0}(C_{g,n})$ 中一族特定的 $C_{g,n}$ 的全纯分歧覆叠 $\Sigma$，其次数为 $N$，并在穿孔附近具有指定行为。$C_{g,n}$ 的分歧覆叠 $\Sigma$ 等同于与所考虑的库仑分支上的点相对应的 Seiberg–Witten 曲线；由于它位于 $T^*_{1,0}(C_{g,n})$ 中，它可以由如下形式的方程定义：
\begin{equation}
\lambda^N+\sum_{i=2}^{N}\phi_i\lambda^{N-i}=0\,,\tag{22.1}
\end{equation}
其中 $\lambda$ 是 $T^*_{1,0}(C_{g,n})$ 上的全纯 Liouville 1-形式，而对每个 $i=2,\dots,N$，$\phi_i$ 是 $C_{g,n}$ 上的全纯 $i$-微分，其在 $C_{g,n}$ 的穿孔处的极点由 $\mathcal{D}$ 规定。

% [Tutor Note: 可与第 20 章的谱曲线方程 (20.19) 对读：那里对 $G=\SU(N)$ 写 $\det(\eta-\phi)=\eta^{\otimes N}+\eta^{\otimes(N-2)}W_2+\eta^{\otimes(N-3)}W_3+W_N=0$，式中同样没有 $\eta^{\otimes(N-1)}$ 项，原因是 $\phi$ 取伴随表示即无迹，$\tr\phi=0$。此处 (22.1) 把 $\eta$ 换成 Liouville 形式 $\lambda$、把 $W_i(\phi)$ 换成 $\phi_i$，正是同一特征方程在 S 类理论语境下的写法（$\lambda|_{\Sigma}$ 即该理论的 Seiberg–Witten 微分 $\lambda_{SW}$）。核对方法：$N=2$ 时 (22.1) 为 $\lambda^2+\phi_2=0$，与 22.2 节的 (22.6) $\lambda^2-\phi_2=0$ 相差 $\phi_2$ 的符号约定（两处 $\phi_2$ 的定义差一符号）。]

% [原文 p.367]
让我们假定所考虑的是理论库仑分支上的一般点，此时覆叠 $\Sigma\to C_{g,n}$ 只有单分支点。在 $C_{g,n}$ 上选取终止于分支点的分支割线后，可以用指标 $i=1,\dots,N$ 标号覆叠的各叶，而分支点与分支割线由对 $\{ij\}$ 标号。

与平坦空间情形一样，M2-膜可以在覆叠 $\Sigma\to C_{g,n}$ 的各叶之间伸展。这样的构型是 BPS 的，当且仅当这些 M2-膜缠绕超凯勒空间 $Q$ 中的特殊拉格朗日圈。在场论极限下，这一条件特别要求 M2-膜是「竖直的」，意即它们在覆叠映射 $\Sigma\to C_{g,n}$ 下投影为 $C_{g,n}$ 上的弦。然而，并非所有这样的构型都必然是 BPS 的。

\subsection{局部 BPS 条件}

设 $(z,x)$ 是 $T^*_{1,0}(C_{g,n})$ 上的局部全纯坐标，其中 $z$ 是 $C_{g,n}$ 的开片 $U$ 上的坐标，$x$ 是 $T^*_{1,0}(C_{g,n})|_U\to U$ 的纤维上的坐标。假定 $U$ 不含分支点与分支割线。用 $z$ 与 $x$ 表示，1-形式 $\lambda$ 读作 $\lambda=x\dd z$，而覆叠 $\Sigma|_U\to U$ 的 $N$ 个叶可以由 $U$ 上的 $N$ 个函数 $x_1(z),\dots,x_N(z)$ 描述。对每个 $i=1,\dots,N$，令
\[
\lambda_i=x_i\dd z
\]
为 Liouville 形式 $\lambda$ 到覆叠第 $i$ 叶的限制，并对任意 $j\ne i$ 令 $\lambda_{ij}=\lambda_i-\lambda_j$。注意 $\lambda_{ij}$ 是 $U$ 上的一个 1-形式，它编码了 $\Sigma|_U\to U$ 的第 $i$ 叶与第 $j$ 叶之间的间距。

\Fig[0.62\linewidth]{figures/fig_22_2.png}{22.2}{$C_{g,n}$ 上来自 M2-膜的弦。}

设 $\gamma_{ij}$ 是 $U$ 中一条有向 $(ij)$-弦的小段，由光滑映射
\[
\gamma_{ij}:\ t\in\,]-1,1[\,\longmapsto \gamma_{ij}(t)\in U
\]
给出。

% [原文 p.368]
这一设定见图 22.2。与平坦空间情形类似，$\gamma_{ij}$ 的张力正比于第 $i$ 叶与第 $j$ 叶之间的间距。不过此时弦的张力沿 $\gamma_{ij}$ 按 $|x_i(\gamma_{ij}(t))-x_j(\gamma_{ij}(t))|$ 变化。这样的弦因而称为变张力弦。沿 $\gamma_{ij}$ 对张力积分得到这段 $(ij)$-弦的质量：
\begin{equation}
M(\gamma_{ij})=\frac{1}{\pi}\int_{\gamma_{ij}}|\lambda_{ij}|\,.\tag{22.2}
\end{equation}
质量 $M(\gamma_{ij})$ 与 $\gamma_{ij}$ 的定向无关。

回忆 $\lambda|_{\Sigma}$ 充当 S 类理论的 Seiberg–Witten 微分，并且沿 Seiberg–Witten 曲线 $\Sigma$ 上的有向圈 $\gamma$ 积分 $\lambda|_{\Sigma}$ 给出 BPS 态的 $4d$ $\mathcal{N}=2$ 中心荷，该态的电荷由 $\gamma$ 在某个辛基下的同调类给出。设 $\gamma$ 是 $\Sigma$ 上的有向圈，使得 $\gamma^{(i)}_{ij}-\gamma^{(j)}_{ij}\subset\gamma$，其中对每个 $k$，$\gamma^{(k)}_{ij}$ 表示 $\gamma_{ij}$ 到覆叠 $\Sigma|_U\to U$ 第 $k$ 叶的提升，而负号表示定向反转。$\gamma^{(i)}_{ij}-\gamma^{(j)}_{ij}$ 对中心荷的贡献为
\begin{equation}
Z(\gamma_{ij}):=\frac{1}{\pi}\bigg(\int_{\gamma^{(i)}_{ij}}\lambda|_{\Sigma}-\int_{\gamma^{(j)}_{ij}}\lambda|_{\Sigma}\bigg)=\frac{1}{\pi}\int_{\gamma_{ij}}\lambda_{ij}\,.\tag{22.3}
\end{equation}
我们略微滥用记号，把 $Z(\gamma_{ij})$ 称作弦段 $\gamma_{ij}$ 的中心荷。注意若 $\gamma^{\mathrm{opp}}_{ij}(t):=\gamma_{ij}(-t)$ 表示定向反转的弦、且 $\gamma_{ji}(t)=\gamma_{ij}(t)$，则 $Z(\gamma^{\mathrm{opp}}_{ij})=Z(\gamma_{ji})$。

BPS 态对应于 $\Sigma$ 中的 1-圈：M2-膜在环绕 $\Sigma$ 的 M5-膜上附着，而这些 1-圈投影为 $C_{g,n}$ 上的弦，或更一般地，投影为弦网。上面的分析表明，这类弦网的质量与中心荷可以在 $C_{g,n}$ 上局部计算。S 类理论中由这样的弦网表示的态是 BPS 的，当且仅当它满足 BPS 条件 $M=|Z|$。等价地，网中每一段 $\gamma_{ij}$ 都必须满足局部 BPS 条件\footnote{本章采用 [155] 的约定，它与第 20 章的约定相差一个因子 $\sqrt{2}$。} $M_{\gamma_{ij}}=|Z_{\gamma_{ij}}|$，其中 $M_{\gamma_{ij}}$ 与 $Z_{\gamma_{ij}}$ 由式 (22.2) 与 (22.3) 定义。更确切地说：

\begin{definition}[22.1]
$C_{g,n}$ 上 $(ij)$-弦的小段 $\gamma_{ij}$ 满足局部 BPS 条件，当且仅当存在某个 $\vartheta\in[0,2\pi[$，使得对任意 $t\in\,]0,1[$，记 $\partial_t$ 为 $\gamma_{ij}$ 在 $\gamma_{ij}(t)\in U$ 处相切的向量，有
\begin{equation}
\lambda_{ij}\cdot\partial_t\in \ee^{\ii\vartheta}\R_+\,.\tag{22.4}
\end{equation}
\end{definition}

当此条件满足时，有 $Z(\gamma_{ij})=\ee^{\ii\vartheta}M(\gamma_{ij})$。

\begin{remark}[22.2]
为使弦网表示一个 BPS 态，式 (22.4) 中的角 $\vartheta$ 必须在构成该网的所有弦段上相同。
\end{remark}

% [原文 p.369]
\begin{remark}[22.3]
注意条件 (22.4) 正是 13.6 节中谱网络的解析定义所用的条件（的一个轻微推广）。实践中，本章所考虑的谱网络大多属于 WKB 型，因为从 S 类理论的角度看这是最自然的框架。[154] 中的构造从这类 WKB 谱网络出发，随后把这一概念抽象为拓扑的概念，引入阿贝尔化与非阿贝尔化两个过程（参见第 15 章）。
\end{remark}

为清楚起见记 $\gamma_{ij}(t)=:z(t)$，则式 (22.4) 可以重新表述为：
\begin{equation}
\big[x_i(z(t))-x_j(z(t))\big]\frac{\dd z}{\dd t}(t)\in \ee^{\ii\vartheta}\R_{>0}\,.\tag{22.5}
\end{equation}

% [Tutor Note: 注意原文在 (22.4) 处写 $\R_+$、在 (22.5) 处写 $\R_{>0}$，二者应作同一理解（正实数集）。可核对的另一点：由 (22.2) 与 (22.3) 可见质量与中心荷共用同一因子 $1/\pi$，故 (22.4) 成立时必有 $Z(\gamma_{ij})=\ee^{\ii\vartheta}M(\gamma_{ij})$，与定义 22.1 之后的断言一致；而若像第 20 章那样把中心荷定义为不带 $1/\pi$ 的 $\oint\lambda$，则两章的中心荷相差因子 $\pi$——这正是脚注 1 所说「与第 20 章相差一个因子」的具体来源之一（另一来源是 $\sqrt2$ 的归一化）。]

\subsection{有限弦网}

为使满足给定角 $\vartheta$ 的局部 BPS 条件的弦网对应于真正的 BPS 态，它还必须是有限的。这一要求来自式 (22.2)：非有限的网将具有无穷大质量。于是我们得到 S 类理论中 BPS 态的如下刻画：

\begin{proposition}[22.4]
S 类理论中的 BPS 态对应于紫外曲线 $C_{\mathrm{UV}}$ 上满足式 (22.4) 的局部 BPS 条件（对某个固定角 $\vartheta$）的弦网。该态的质量与中心荷由构成该网的各弦段贡献之和得到（参见式 (22.2) 与 (22.3)）。特别地，中心荷的辐角为 $\vartheta$。
\end{proposition}

有限 BPS 弦的例子及其 M 理论来源见图 22.3。

\Fig[0.85\linewidth]{figures/fig_22_3.png}{22.3}{作为有限 M2-膜与作为有限弦的 BPS 态。}

% [原文 p.370]
理论上，这为计算与分析 S 类理论的 BPS 态提供了一个框架，尽管实践中并没有实现这一方法的直接算法。解释如何做到这一点将是接下来几节的重点。

\begin{remark}[22.5]
这种处理 $4d$ $\mathcal{N}=2$ 理论中 BPS 态的方法，最早是在由弦论构造出的规范理论语境中开创的 [233, 65, 328]。
\end{remark}

\section{$A_1$ 型中的 BPS 叶状结构}

\subsection{局部描述}

在本节中，我们专注于 $A_1$ 型的 S 类理论。在这些理论中，Seiberg–Witten 曲线 $\Sigma$ 是紫外曲线 $C_{g,n}$ 的 2 次分歧覆叠，由方程
\begin{equation}
\lambda^2-\phi_2=0\,,\tag{22.6}
\end{equation}
确定，其中 $\lambda$ 是 $T^*C_{g,n}$ 上的全纯 Liouville 1-形式，$\phi_2$ 是 $C_{g,n}$ 上的二次微分。对任意给定的角 $\vartheta\in[0,2\pi[$，式 (22.4) 中的局部 BPS 条件定义 $C_{g,n}$ 的一个叶状结构，下面我们描述它。该叶状结构的所有叶都是 $(12)$-弦（或 $(21)$-弦），本节中我们常把它们简称为 BPS 弦。

\begin{remark}[22.6]
这些由二次微分定义的黎曼曲面叶状结构，与 8.2 节（另见 16.6 节）中引入并研究的叶状结构相同；我们这里从 S 类理论的角度重新审视它们。
\end{remark}

\subsubsection*{正则点}

在 $C_{g,n}$ 中正则点的充分小开邻域 $U$ 内，函数 $x_{1,2}(z)$ 可以用常数逼近。此时 BPS 微分方程 (22.4) 约化为
\[
\frac{\dd z}{\dd t}(t)=\alpha,
\]
其中 $\alpha$ 是辐角固定的复常数。该微分方程在 $U$ 内的积分曲线是斜率为 $\arg(\alpha)$ 的平行直线。

\subsubsection*{分支点}

考虑一个（单）分支点 $B=B_{12}\in C_{g,n}$。设 $z$ 是 $B$ 附近的全纯坐标，使得 $B$ 位于 $z=0$，且在 $B$ 附近（对平方根 $\sqrt{z}$ 的某个分支选取）有 $x_1(z)\sim\sqrt{z}$、$x_2(z)\sim-\sqrt{z}$。

% [原文 p.371]
令 $U=|z|<\epsilon$（$\epsilon$ 适当小）。我们选 $z\in\R_{<0}\cap U$ 为分支割线。在此记号下，式 (22.4) 变为
\[
2\sqrt{z}\,\frac{\dd z}{\dd t}(t)=\frac{4}{3}\frac{\dd}{\dd t}z^{3/2}(t)\in \ee^{\ii\vartheta}\R_{>0}\,.
\]
该微分方程在 $U$ 内的积分曲线中有三条 BPS 弦，每条都起于或止于该分支点。这些弦可由 $z_k(t)\propto t^{2/3}\ee^{\ii\phi}$ 描述，其中 $\phi=(2\vartheta+4k\pi)/3$，$k=0,1,2$。由于分支割线的存在，$z_0$ 表示一条从 $B$ 出发的 $(12)$-曲线，而 $z_1$ 与 $z_2$ 对应于终止于 $B$ 的 $(12)$-曲线。不计定向可以看出，在每个分支点 $B$ 处都有三条弦从 $B$ 发出。$U$ 中其余的积分曲线在远离 $B$ 处趋于 $z_1$、$z_2$ 或 $z_3$。这一局部构型见图 22.4。

\Fig[0.4\linewidth]{figures/fig_22_4.png}{22.4}{分支点 $B$ 附近局部 BPS 方程 (22.4) 的积分曲线。所有带定向的实线都是 $(12)$-弦，虚线是分支割线。}

\subsubsection*{正则穿孔}

现在我们转向式 (22.4) 在 $C_{g,n}$ 的正则穿孔附近的积分曲线描述。设 $z$ 是穿孔邻域内的全纯坐标，使穿孔位于 $z=0$。设 $m\in\C$ 是与该穿孔相伴的质量参数，在 $A_1$ 型情形它必然是单的（或等价地，满的）。二次微分 $\phi_2$ 具有系数为 $m^2$ 的二阶极点，而式 (22.6) 的解满足 $x_\pm(z)\dd z\sim\pm m\dd z/z$。于是式 (22.4) 的 BPS 微分方程读作
\[
\frac{2m}{z}\,\frac{\dd z}{\dd t}\in \ee^{\ii\vartheta}\R_{>0}\,.
\]
穿孔附近的积分曲线可以参数化为
\[
z(t)=C\exp(\alpha t),
\]

% [原文 p.372]
其中 $C$ 与 $\alpha$ 是复常数，且 $\alpha=\ee^{\ii\vartheta}/m$。使 $\operatorname{Re}(\alpha)=0$ 的 $\vartheta$ 值很特殊：此时积分曲线绕着穿孔转圈。对其他的 $\vartheta$ 值，积分曲线或者当 $t\to-\infty$ 时从穿孔发出，或者当 $t\to\infty$ 时收敛于它，从而描出对数螺线。具体地，设 $\alpha=\alpha_1+\ii\alpha_2$，则
\[
z(t)=C\ee^{\alpha_1 t}\big(\cos(\alpha_2 t)+\ii\sin(\alpha_2 t)\big).
\]
这类积分曲线见图 22.5。

\Fig[0.9\linewidth]{figures/fig_22_5.png}{22.5}{正则穿孔邻域内局部 BPS 方程的积分曲线。}

\subsubsection*{非正则穿孔}

在 21.4.4 节中我们看到，非正则穿孔的特征是定义分歧覆叠 $\Sigma\to C_{g,n}$ 的全纯微分在穿孔处具有更高阶的极点。这类穿孔典型地出现在 $F\le3$ 个味的 $SU(2)$ SYM 理论语境中，此时二次微分 $\phi_2$ 中出现 3 阶与 4 阶极点。设 $z$ 是非正则穿孔邻域内的全纯坐标，穿孔位于 $z=0$，并设 $k+2$ 表示二次微分 $\phi_2$ 的阶，其中 $k$ 是正整数：
\[
\phi_2(z)\sim\frac{c}{z^{k+2}}\,.
\]
在穿孔附近，$\Sigma$ 的各分支行为如 $x_\pm(z)\sim\pm\sqrt{c}\,z^{-k/2-1}$，因此式 (22.4) 变为
\[
\frac{2\sqrt{c}}{z^{k/2+1}}\,\frac{\dd z}{\dd t}\in \ee^{\ii\vartheta}\R_{>0}\,.
\]
该方程的积分曲线可以参数化为
\[
z(t)=\Big(\frac{1}{\alpha t+C}\Big)^{2/k},
\]

% [原文 p.373]
其中 $C$ 是复常数，$\alpha=-k\ee^{\ii\vartheta}/(4\sqrt{c})$。当 $t\to\infty$ 时，这些积分曲线沿 $k$ 个不同的渐近方向逼近穿孔。$k=1$ 与 $k=4$ 的例子见图 22.6。

% [Tutor Note: 可直接验证非正则穿孔处的参数化解。取 $z(t)=(\alpha t+C)^{-2/k}$，则 $\dd z/\dd t=-(2\alpha/k)(\alpha t+C)^{-2/k-1}$，而 $z^{-k/2-1}=(\alpha t+C)^{(2/k)(k/2+1)}=(\alpha t+C)^{1+2/k}$；两者相乘得 $2\sqrt c\,z^{-k/2-1}\dd z/\dd t=2\sqrt c\,(-2\alpha/k)=-4\sqrt c\,\alpha/k=\ee^{\ii\vartheta}$（代入 $\alpha=-k\ee^{\ii\vartheta}/(4\sqrt c)$），恰落在 $\ee^{\ii\vartheta}\R_{>0}$ 中。同理可核正则穿孔：$z=C\ee^{\alpha t}$ 时 $(2m/z)\dd z/\dd t=2m\alpha=2\ee^{\ii\vartheta}$（代入 $\alpha=\ee^{\ii\vartheta}/m$），也落在 $\ee^{\ii\vartheta}\R_{>0}$ 中。]

\begin{remark}[22.7]
局部 BPS 方程的积分曲线在非正则穿孔附近的行为，与 4.3 节引入的纤毛曲面概念相关。对二次微分 $\phi_2$ 的阶为 $k=2$ 的非正则穿孔，我们可以挖去穿孔周围的一个小圆盘，并标记该圆盘边界与「积分曲线收敛到穿孔所沿的渐近方向」之间的 $k$ 个交点，从而引入一个带 $k$ 根纤毛的边界分支。我们把这一纤毛曲面记作 $C_{\mathrm{UV}}$。
\end{remark}

% [Tutor Note: 原文此处「$\phi_2$ 的阶为 $k=2$」疑为笔误，应为「阶为 $k+2$」（$k$ 一般）。依据有三:(i) 本节定义中明确写「$k+2$ 表示二次微分 $\phi_2$ 的阶」；(ii) 后文 p.380 说「在该穿孔处 $\phi_2$ 有 $k+2$ 阶极点，$k\ge3$」，并得到 $(k+2)$ 边形的 WKB 三角剖分；(iii) 图 22.6 标注 $k=4$ 时恰有四根纤毛。故「$k$ 个交点、$k$ 根纤毛」应对一般的 $k$。此处照原文译出，读者按 $k+2$ 理解。]

\Fig[0.98\linewidth]{figures/fig_22_6.png}{22.6}{非正则穿孔邻域内局部 BPS 方程的积分曲线。}

以上完成了我们对局部 BPS 方程 (22.4) 的积分曲线的局部描述。

\subsection{WKB 三角剖分}

由式 (22.4) 定义的叶状结构中的一条给定弦必属于以下之一：
\begin{itemize}
\item 一般弦，起于一个穿孔并止于一个穿孔（可以是同一个）；
\item 分离弦，或起于穿孔止于分支点，或反之；
\item 有限弦，起于一个分支点止于一个分支点，或构成闭合圈。这类有限弦及其 M 理论来源见图 22.3；
\item 发散弦，它不闭合，且至少在一个方向上没有极限。
\end{itemize}

\begin{remark}[22.8]
当 $C_{\mathrm{UV}}$ 至少有一个穿孔或纤毛时，在不存在有限弦的前提下，发散弦不可能存在 [336]。使不存在有限弦的 $\vartheta$ 值称为一般的。
\end{remark}

因此，对一般的 $\vartheta$ 值，BPS 弦或者是一般弦或者是分离弦，而一般的 BPS 弦构成以一族

% [原文 p.374]
分离 BPS 弦之并为边界的 1-参数族。这些 1-参数族有两种可能的拓扑，取决于分离弦是与两个分支点还是与单个分支点相伴。这两种拓扑见图 22.7，其中一般的 BPS 弦用黑色实线表示，分离 BPS 弦用红色虚线表示，穿孔用黑点表示，分支点用红色星号表示。

\Fig[0.8\linewidth]{figures/fig_22_7.png}{22.7}{一般 BPS 弦的 1-参数族。}

从一般 BPS 弦的每个 1-参数族中选出一个代表，就得到 $C_{\mathrm{UV}}$ 在定义 4.16 意义下的一个三角剖分，称为对应于角 $\vartheta$ 的 $C_{\mathrm{UV}}$ 的 WKB 三角剖分。我们下面按原始文献 [151, 155] 简要解释这一术语背后的理由，细节请参见原文。

\begin{remark}[22.9]
这里定义的 WKB 三角剖分与 8.3 节中由二次微分的临界图引入的三角剖分一致，区别只在于：如我们马上会看到的，WKB 三角剖分比单纯的三角剖分编码了略多的数据；它是一个带装饰的三角剖分，这一概念在 [155] 中引入，参照的是 [129] 中的带装饰局部系统\footnote{严格说来，在 [155] 中起作用的不只是装饰本身，而是标架化。}，参见 10.2 节。
\end{remark}

\subsection{BPS 态计数}

\subsubsection*{BPS 指标}

为澄清「BPS 计数」的含义，注意一般不能指望 $4d$ $\mathcal{N}=2$ 理论库仑分支 $\mathcal{C}$ 中给定点 $u\in\mathcal{C}$ 处的 BPS 态数目在 $u$ 变化时保持不变，即使没有穿过边缘稳定性墙。

% [原文 p.375]
如第 20 章所讨论，$\mathcal{C}$ 中一般点处低能激发的电荷构成 $\mathcal{C}$ 上的格局部系统 $\Gamma=(\Gamma_u)_{u\in\mathcal{C}}$，其秩为 $2r+f$，其中 $r$ 表示理论的秩（即库仑分支的维数），$f$ 是 $\U(1)$ 味对称性因子的个数。这个局部系统带有反对称的 Dirac–Schwinger–Zwanziger 配对 $\langle\cdot,\cdot\rangle$，它在味荷上平凡。局部系统 $\Gamma$ 在 $\mathcal{C}$ 中某些 BPS 粒子变为无质量的复余维 1 轨迹周围表现出单值变换。此外，$\Gamma$ 上有一个复值线性函数 $Z(u)$，编码中心荷，即 $Z(u)\in\Gamma_u^*\otimes_{\Z}\C$，它随 $u$ 全纯变化。函数 $Z(u)$ 通过在 $u$ 处的 Seiberg–Witten 曲线 $\Sigma_u$ 上的 1-圈上积分 Seiberg–Witten 微分 $\lambda$ 来计算。具体地，对每个 $\gamma\in\Gamma_u$ 有
\[
Z_\gamma(u):=Z(u)(\gamma)=\oint_{\widehat{\gamma}}\lambda_u,
\]
其中 $\lambda_u$ 是 $u\in\mathcal{C}$ 处的 Seiberg–Witten 微分，$\widehat{\gamma}$ 是 $\Sigma_u$ 上同调类由 $\gamma\in\Gamma_u$ 给出的 1-圈。此外，每个元素 $\gamma\in\Gamma_u$ 都决定 $u\in\mathcal{C}$ 处 1-粒子态 Hilbert 空间的一个子空间 $\mathcal{H}_{\gamma,\mathrm{BPS}}(u)$，它由理论中电、磁与味荷均为 $\gamma$、且满足 BPS 界 $M\ge|Z_\gamma(u)|$ 的所有单粒子态组成。

BPS 态不能以直接的方式计数：人们不能指望在 $\mathcal{C}$ 中边缘稳定性墙的补集的每个连通分支上都有一个良定义的整数，原因是（半）BPS 态在 $u$ 变化时可以成对重新组合。反过来，超对称代数的非 BPS 表示也可以分裂为成对的 BPS 态。为此，改用超对称的「受保护」指标，类似于 [364, 74] 中定义的受保护指标。

\begin{proposition}[22.10]
对具有 $\mathcal{N}=2$ 超对称的四维理论中的 BPS 态，相应的指标称为第二螺旋度超迹，写作
\begin{equation}
\Omega(\gamma;u)=-\frac{1}{2}\operatorname{Tr}_{\mathcal{H}^\gamma_{1,\mathrm{BPS}}(u)}(-1)^{2J_3}(2J_3)^2\,,\tag{22.7}
\end{equation}
其中 $J_3$ 是有质量小群的转动子群的生成元。
\end{proposition}

这个指标在计入上述含混性的同时「计数」BPS 粒子，意即它在 $\mathcal{C}$ 中边缘稳定性墙补集的各连通分支上保持常数。如第 20 章所讨论，BPS 谱在穿过边缘稳定性墙时可以改变；相应地，$\Omega(\gamma;u)$ 可以在这些墙处跳跃。可以证明，每个 BPS 超多重态对 $\Omega(\gamma;u)$ 贡献因子 $1$，而每个 BPS 矢量多重态贡献因子 $-2$。

% [原文 p.376]
\subsubsection*{Kontsevich–Soibelman 穿墙公式}

在 [237] 中，Kontsevich 与 Soibelman 为三维 Calabi–Yau 流形的 Donaldson–Thomas 不变量 $\widehat{\Omega}(\gamma;u)$ 导出了一个穿墙公式。尽管这一语境看起来与 $4d$ $\mathcal{N}=2$ 理论中的 BPS 计数无关，他们的穿墙公式却适用于第二螺旋度超迹 $\Omega(\gamma;u)$。这一联系在 [151] 中通过考察在圆周上紧化的 $4d$ $\mathcal{N}=2$ 理论（即在 $\R^3\times S^1_R$ 上，$R$ 表示 $S^1$ 的半径）而得到证明。在大半径与低能下，有效理论是一个三维 $\mathcal{N}=4$ sigma 模型，映入可积系统 $\mathbb{M}\to\mathcal{C}$，后者刻画原 $4d$ $\mathcal{N}=2$ 理论的低能动力学 [322]。由于 Seiberg–Witten 可积系统的底对应于 $4d$ $\mathcal{N}=2$ 理论的库仑分支 $\mathcal{C}$，其纤维（一般地是 $2r$ 维环面）由「电」($\theta_e^I$) 与「磁」($\theta_m^I$) Wilson 圈的期望值参数化：
\[
\theta_e^I=\oint_{S^1}A^I_4\,\dd x^4\,,\qquad \theta_{m,I}=\oint_{S^1}(A_{D,4})_I\,\dd x^4,
\]
其中 $A^I_4$ 是第 $I$ 个低能 $4d$ 阿贝尔规范场沿 $x^4$ 的分量，而 $A_{D,4}$ 是三维中的对偶光子。三维中的 $\mathcal{N}=4$ 超对称要求 sigma 模型的目标流形是超凯勒的。

Kontsevich 与 Soibelman 的穿墙公式（粗略地）如下导出。格 $\Gamma_u$ 定义复分裂环面 $\widetilde{T}_u=\Gamma_u^*\otimes_{\Z}\C^\times$，而每个 $\gamma\in\Gamma_u$ 自然地定义 $\widetilde{T}$ 上的函数 $X_\gamma$。这些函数满足关系 $X_{\gamma_1}X_{\gamma_2}=X_{\gamma_1+\gamma_2}$，且 $\Gamma_u$ 的任一组基定义 $\widetilde{T}$ 上的坐标系。此外，如 5.5 节那样，DSZ 配对 $\langle\cdot,\cdot\rangle$ 在 $\widetilde{T}_u$ 上诱导一个全纯辛形式。设 $\sigma$ 是 $(-1)^{\langle\gamma_1,\gamma_2\rangle}$ 的二次细化，即
\[
\sigma(\gamma_1)\sigma(\gamma_2)=(-1)^{\langle\gamma_1,\gamma_2\rangle}\sigma(\gamma_1+\gamma_2),
\]
并设 $e_\gamma$ 是由「哈密顿量」$\sigma(\gamma)X_\gamma$ 生成的 $\widetilde{T}_u$ 的辛同胚。所有对 $\gamma\in\Gamma_u$ 的 $e_\gamma$ 生成一个李代数，其关系为
\[
[e_{\gamma_1},e_{\gamma_2}]=(-1)^{\langle\gamma_1,\gamma_2\rangle}\langle\gamma_1,\gamma_2\rangle e_{\gamma_1+\gamma_2}\,.
\]
现在令
\[
K_\gamma=\exp\sum_{n=1}^{\infty}\frac{1}{n^2}e_{n\gamma}\,.
\]
这个变换作用在函数 $X_{\gamma'}$ 上为
\begin{equation}
K_\gamma:\ X_{\gamma'}\longrightarrow X_{\gamma'}\big(1-\sigma(\gamma)X_\gamma\big)^{\langle\gamma',\gamma\rangle}.\tag{22.8}
\end{equation}

\begin{remark}[22.11]
这些是簇变换；参见 5.5 节与 [237]。
\end{remark}

% [Tutor Note: 可用第 5 章的变异公式核对 (22.8)。把 $X_\gamma$ 看作簇坐标 $x_i$、(22.8) 即 $X_{\gamma'}\mapsto X_{\gamma'}(1-X_\gamma)^{\langle\gamma',\gamma\rangle}$；与式 (5.1) 的 X 型变异 $\mu_i(x_j)=x_j(1+x_i^{\sgn\epsilon_{ij}})^{\epsilon_{ij}}$ 相比，只差一个符号约定：取 $X_\gamma$ 为「$-x_i$」并把指数 $\langle\gamma',\gamma\rangle$ 对应到 $-\epsilon_{ij}$，两式一致。这也解释了为何 (22.10) 的乘积在穿墙时不变：它就是簇代数里五边形关系 $K_1K_2K_1=K_2K_1K_2$ 型的关系式（对照命题 5.10 与例 22.14 的 $K_{1,0}K_{0,1}=K_{0,1}K_{1,1}K_{1,0}$)。]

设 $u_0\in\mathcal{C}$ 是四维库仑分支 $\mathcal{C}$ 中边缘稳定性墙上的一个一般点，$u_+$ 与 $u_-$ 是位于墙两侧靠近 $u_0$ 的点。

% [原文 p.377]
对每个 $u\in\mathcal{C}$ 以及每个电荷为 $\gamma$ 的 BPS 态，令 $l_{\gamma,u}\subset\C$ 表示半直线
\begin{equation}
l_{\gamma,u}=\{\zeta\cdot Z_\gamma(u)\mid \zeta\in\R_-\}\,.\tag{22.9}
\end{equation}
设 $Z_{\gamma_1}$ 与 $Z_{\gamma_2}$ 是在 $u_0$ 所在的边缘稳定性墙上对齐的中心荷。在支撑性质（support property）的假设下，在墙上变为无质量的粒子（相对于反粒子）的电荷属于集合
\[
\mathcal{ML}=\{n\gamma_1+m\gamma_2\mid n,m>0\}.
\]
注意当穿过墙时，对 $\gamma\in\mathcal{ML}$，射线 $l_{\gamma,u}$ 在 $\C$ 中的次序会改变。最后，对远离墙的 $u$，令
\[
\overset{\curvearrowright}{\prod_{\gamma\in\mathcal{ML}}}K_\gamma^{\Omega(\gamma;u_+)}
\]
表示按 $l_{\gamma,u}$ 在 $\C$ 中的顺时针定向排序的 $K_\gamma^{\Omega(\gamma;u_+)}$ 之积。

\begin{proposition}[22.12（Kontsevich–Soibelman 穿墙公式）]
\[
\overset{\curvearrowright}{\prod_{\gamma\in\mathcal{ML}}}K_\gamma^{\Omega(\gamma;u_+)}=\overset{\curvearrowright}{\prod_{\gamma\in\mathcal{ML}}}K_\gamma^{\Omega(\gamma;u_-)}.\tag{22.10}
\]
\end{proposition}

在原始的穿墙公式中，Donaldson–Thomas 不变量 $\widehat{\Omega}(\gamma;u)$ 取代式 (22.10) 中的第二螺旋度超迹 $\Omega(\gamma;u)$。然而，Gaiotto、Moore 与 Neitzke 证明了该公式也适用于 $4d$ $\mathcal{N}=2$ 理论中的 BPS 态，并且它实际上可以由有效 $3d$ $\mathcal{N}=4$ sigma 模型目标空间 $\mathbb{M}$ 上超凯勒度量的光滑性导出——该 sigma 模型是把四维理论在圆周 $S^1_R$ 上紧化得到的。朴素的维数约化给出 $\mathbb{M}$ 上的超凯勒度量 $g^{\mathrm{sf}}_{\mathbb{M}}$，它在 $\mathcal{C}$ 中的奇异轨迹上带有奇点。$\mathbb{M}$ 上真正的超凯勒度量 $g_{\mathbb{M}}(R)$ 只接收 $4d$ BPS 粒子的贡献，这些贡献把这些奇点抹平。两个度量之间的关系是：当 $R\to\infty$ 时 $g_{\mathbb{M}}(R)\sim g^{\mathrm{sf}}_{\mathbb{M}}$。当穿过 $\mathcal{C}$ 中的一面边缘稳定性墙时，BPS 谱改变；然而可以预期度量 $g_{\mathbb{M}}$ 光滑地变化。这对穿墙施加了严格的限制，特别是它必须遵循 Kontsevich–Soibelman 穿墙公式（式 (22.10)）。

\begin{example}[22.13]
Kontsevich–Soibelman 穿墙公式的一个特例是
\[
K_{2,-1}K_{0,1}=(K_{0,1}K_{2,1}K_{4,1}\dots)K_{2,0}^{-2}(\dots K_{6,-1}K_{4,-1}K_{2,-1}).
\]
与 20.1.5 节的讨论相比较可以看出，这个方程编码了纯 $SU(2)$ SYM 中观察到的穿墙。
\end{example}

\begin{example}[22.14]
穿墙公式的某些实例比一般情形更简单，因为式 (22.10) 左右两边只涉及有限多项，例如：
\end{example}

% [原文 p.378]
\[
K_{1,0}K_{0,1}=K_{0,1}K_{1,1}K_{1,0}\,.
\]
这个公式编码了一个简单 Argyres–Douglas 理论中的穿墙（例如参见 [328]）。

从超凯勒度量的扭量构造（twistorial construction）[201, 200] 的角度可以理解 $4d$ BPS 态如何对「朴素」超凯勒度量 $g^{\mathrm{sf}}_{\mathbb{M}}$ 给出修正。作为超凯勒流形，有效 $3d$ $\mathcal{N}=4$ sigma 模型的目标空间 $\mathbb{M}$ 容许一个 $\C\PP^1$-值的复结构。编码 $\mathbb{M}$ 上复结构的空间 $\C\PP^1$ 称为与 $\mathbb{M}$ 相伴的扭量球面，而与超凯勒流形 $\mathbb{M}$ 相伴的扭量空间是乘积 $\mathbb{M}\times\C\PP^1$，它自然地带有复流形结构。设 $\zeta$ 是扭量球面 $\C\PP^1_\zeta$ 上的标准坐标。

拓扑上，$\mathbb{M}$ 是 $\mathcal{C}$ 上的环面纤维化，即它可以局部由 $(u,\theta)\in\mathbb{M}$ 参数化。$\mathbb{M}$ 上的超凯勒度量由一个映射
\[
\mathcal{X}^{\mathrm{RH}}:\ \mathbb{M}\times\C^\times_\zeta\longrightarrow \Gamma^*\times\C^\times
\]
构造。$\mathcal{X}^{\mathrm{RH}}$ 与电荷格中某个元素 $\gamma\in\Gamma$ 的缩并记作 $\mathcal{X}^{\mathrm{RH}}_\gamma(u,\theta,\zeta)$。有 $\mathcal{X}^{\mathrm{RH}}_\gamma\mathcal{X}^{\mathrm{RH}}_{\gamma'}=\mathcal{X}^{\mathrm{RH}}_{\gamma+\gamma'}$。$\Gamma$ 上的 DSZ 配对定义一个泊松结构：
\[
\big\{\mathcal{X}^{\mathrm{RH}}_\gamma,\mathcal{X}^{\mathrm{RH}}_{\gamma'}\big\}=\langle\gamma,\gamma'\rangle\mathcal{X}^{\mathrm{RH}}_{\gamma+\gamma'}\,.
\]
此外，对每个使 $\Omega(\gamma;u)\ne0$ 的 $\gamma\in\Gamma_u$，要求 $\mathcal{X}^{\mathrm{RH}}_\gamma$ 在扭量球面 $\C\PP^1_\zeta$ 上、除式 (22.9) 定义的半直线 $l_{\gamma,u}$ 之外是全纯的。当 $\zeta$ 穿过 $\C\PP^1_\zeta$ 中的直线 $l_{\gamma,u}$ 时，函数 $\mathcal{X}^{\mathrm{RH}}_\gamma$ 不连续，其间断由式 (22.8) 那样的 Kontsevich–Soibelman 变换给出。于是可以在 $\mathbb{M}$ 上构造一个全纯辛形式与一个超凯勒度量，使得 $\mathcal{X}_\gamma(u,\theta,\zeta)$ 充当全纯 Darboux 坐标。

对应于半平坦度量 $g^{\mathrm{sf}}$ 的函数 $\mathcal{X}^{\mathrm{sf}}_\gamma(u,\theta,\zeta)$ 可以显式表为
\[
\mathcal{X}^{\mathrm{sf}}_\gamma(u,\theta,\zeta)=\exp\Big(\pi R\zeta^{-1}Z_\gamma+\ii\theta_\gamma+\pi R\zeta\overline{Z}_\gamma\Big),
\]
其中 $\theta_\gamma:\Gamma^*\otimes\R/2\pi\Z\to\R/2\pi\Z$ 是与 $\Gamma$ 上坐标对偶的坐标。对应于 $\mathbb{M}$ 上真正超凯勒度量的函数 $\mathcal{X}^{\mathrm{RH}}_\gamma$ 则通过计入来自 $4d$ BPS 粒子的修正而得到。结果是，当 $\zeta\to0,\infty$ 以及当 $R\to\infty$ 时，$\mathcal{X}^{\mathrm{RH}}_\gamma$ 都等价于 $\mathcal{X}^{\mathrm{sf}}_\gamma$。$\mathcal{X}^{\mathrm{RH}}_\gamma$ 中的上标「RH」表明这些函数是一个无穷维 Riemann–Hilbert 问题的解。

% [原文 p.379]
\subsection{$A_1$ 型的 S 类理论}

在 $A_1$ 型 S 类理论这一特殊情形中，模空间 $\mathbb{M}$ 对应于紫外曲线 $C_{\mathrm{UV}}$ 上 $SU(2)$ Higgs 丛的 Hitchin 模空间。给定 Hitchin 方程的一个解 $(A,\phi)$，对任意 $\zeta\in\C^\times$ 可以如下构造 $C_{\mathrm{UV}}$ 上的 $\SL_2(\C)$-联络 $\mathcal{A}$：
\[
\mathcal{A}=\frac{R}{\zeta}\phi+A+R\zeta\,\overline{\phi}.
\]
$(A,\phi)$ 解 Hitchin 方程这一条件保证 $\mathcal{A}$ 是平坦的 $\SL_2(\C)$-联络。反之，若 $\mathcal{A}$ 对所有 $\zeta\in\C^\times$ 都平坦，则 $(A,\phi)$ 满足 Hitchin 方程。这一对应促使我们研究 $C_{g,n}$ 上在穿孔处具有指定单值的平坦 $\SL_2(\C)$-联络的模空间。

在这个模空间上可以引入一组值得注意的坐标：Fock–Goncharov 坐标（参见 10.2 节以及 9.2 节末尾的讨论）。[155] 中使用的坐标与 [129] 中定义的坐标密切相关，只有一些小的调整。我们已在 22.2.2 节解释了 $C_{\mathrm{UV}}$ 上变张力弦的局部 BPS 条件如何定义 $C_{\mathrm{UV}}$ 的一个三角剖分。事实上，如注 22.9 所暗示，WKB 三角剖分比单纯的三角剖分编码略多的数据：它是一个带装饰的三角剖分，意即 $C_{\mathrm{UV}}$ 的一个三角剖分连同平坦性方程 $(\dd+\mathcal{A})s=0$ 在 $C_{\mathrm{UV}}$ 的穿孔或纤毛附近的一个解 $s$（定义到相差一个标量）。$C_{\mathrm{UV}}$ 上的 WKB 三角剖分在每个穿孔或纤毛处自然地配备一个典范装饰。这个装饰来自微分方程 $(\dd+\mathcal{A})s=0$ 当 $\zeta\to0$ 时的 WKB 分析（参见第 16 章），从而为「WKB 三角剖分」这一术语提供了依据。

$C_{\mathrm{UV}}$ 的每个 WKB 三角剖分都定义 Fock–Goncharov 坐标
\[
\mathcal{X}^\vartheta_\gamma:\ \mathbb{M}\times\C^\times\longrightarrow\C,
\]
其中 $\vartheta$ 是决定该 WKB 三角剖分的角。函数 $\mathcal{X}^{\mathrm{RH}}_\gamma$ 然后通过把 $\vartheta$ 取为 $\arg(\zeta)$ 而与坐标 $\mathcal{X}^\vartheta_\gamma$ 等同。

当 $\vartheta$ 变化时，$C_{\mathrm{UV}}$ 的 WKB 三角剖分一般保持不变，但当 $\vartheta$ 穿过对应于 BPS 态角度 $\vartheta=\arg(Z_\gamma)$ 的特定值时会发生离散跳跃。联系不同三角剖分的变换可以取如下形式：
\begin{itemize}
\item 翻转（flip），已在本书第 II 与第 III 部分讨论过，它由 BPS 超多重态引起；
\item 抛接（juggle），与极限三角剖分的存在相关，在物理上对应于 BPS 矢量多重态；
\item pop，对应于 $C_{\mathrm{UV}}$ 的某个穿孔或纤毛处平坦截面的改变，不对应于 BPS 态。
\end{itemize}

\begin{remark}[22.15]
翻转与 pop 正是 16.7 节中结合 Voros 符号讨论的那些操作。
\end{remark}

% [原文 p.380]
当三角剖分被翻转时，Fock–Goncharov 坐标经历簇变换。这些簇变换是式 (22.8) 所述 Kontsevich–Soibelman 变换的特例。函数 $\mathcal{X}^{\mathrm{RH}}_\gamma$ 在抛接与 pop 下的变换也可以显式算出，它们对应于 Kontsevich–Soibelman 辛同胚对 $\mathcal{X}^{\mathrm{RH}}_\gamma$ 的作用。

让 $\vartheta$ 从一般值 $\vartheta$ 变到 $\vartheta+\pi$，可以看到 $C_{\mathrm{UV}}$ 在 $\vartheta+\pi$ 处的 WKB 三角剖分在拓扑上与 $\vartheta$ 处的 WKB 三角剖分相同。两个带装饰的三角剖分通过在每个穿孔或纤毛处做一次 pop 相联系。这一变换称为全 pop（omnipop）。相应的变换由
\[
\mathbf{S}=\overset{\curvearrowleft}{\prod_{\{\gamma\mid \vartheta<\arg(-Z_\gamma(u))<\vartheta+\pi\}}}K_\gamma^{\Omega(\gamma;u)}
\]
给出，它联系函数 $\mathcal{X}^\vartheta$ 与 $\mathcal{X}^{\vartheta+\pi}$。这称为谱生成元，因为完全知道 $\mathbf{S}$ 就确定了 BPS 重数 $\Omega(\gamma;u)$。谱生成元 $\mathbf{S}$ 的计算在 [155, 第 11 节] 中完成，并给出了大量例子。

最简单的例子之一是秩 1 的 Argyres–Douglas 理论，其紫外曲线是一个带单个非正则穿孔的球面。在该穿孔处，二次微分 $\phi_2$ 有一个 $k+2$ 阶极点，其中 $k\ge3$。回顾注 22.7，我们得到 $k+2$ 边形的 WKB 三角剖分。$k=3,4$ 的情形特别简单。$k=3$ 时没有 BPS 态，而 $k=4$ 时有一个 BPS 态，且没有穿墙。$k=5$ 的情形对应于五边形的三角剖分，而 5.5 节讨论的五边形恒等式对应于例 22.14 中那个简单的穿墙公式：
\[
K_{1,0}K_{0,1}=K_{0,1}K_{1,1}K_{1,0}\,.
\]
带 $F\le4$ 个味的 $SU(2)$ 超 Yang–Mills 理论的情形在 [155, 第 10 节] 中有广泛讨论。对 $F=0$，紫外曲线被改写为一个环带，其两个边界分支各有一根纤毛，由此得到例 22.13 的穿墙公式。对 $F=1$，紫外曲线是一个环带，其中一个边界分支有一根纤毛，另一个有两根纤毛。对 $F=2$，有两种标准实现：一种以环带为紫外曲线，两个边界分支各两根纤毛；另一种以圆盘为紫外曲线，其边界分支有一根纤毛，内部有两个穿孔。$F=3$ 与 $F=4$ 的情形类似处理；例如 $F=4$ 情形的「平衡」实现对应于以带四个穿孔的球面为紫外曲线。

\section{谱网络}

谱网络是前两节所述方法的深远推广，它使得研究一般 S 类理论中的 BPS 谱成为可能。

% [原文 p.381]
谱网络与 $A_1$ 型理论情形中引入的变张力弦网非常相似，主要区别在于解释：谱网络是相对于 S 类理论中与 $C_{\mathrm{UV}}$ 相对应的某些面缺陷来定义的。相应地，网的每条边带有一个指标 $\mu$，称为孤子指标。

在这样的面缺陷存在时，可以定义带标架 BPS 态的概念，即与某些超对称界面相伴的 BPS 态。带标架 BPS 态的谱最终是超定的，这又使得可以计算网每条边的孤子指标。此外，[152] 中导出的 Kontsevich–Soibelman 穿墙公式的推广（它计入 $2d$ 与 $4d$ BPS 态之间的耦合）使得可以直接从 $C_{\mathrm{UV}}$ 上有限 BPS 网的拓扑计算 $4d$ BPS 简并度。

我们遵循开创性文献 [154]；特别地，我们限于 $N\ge2$ 的 $A_{N-1}$ 型 S 类理论。谱网络在一般 S 类理论语境中的考察见 [249]。

\subsection{BPS 态、孤子与带标架 BPS 态}

\subsubsection*{BPS 态}

回忆 22.1.3 节给出的 $4d$ BPS 态的几何描述：给定角 $\vartheta$，电荷为 $\gamma$、中心荷为 $Z_\gamma$ 且满足 $\arg(Z_\gamma)=\vartheta$ 的 BPS 态对应于紫外曲线 $C_{\mathrm{UV}}$ 上的一张有限弦网。网的每条边\footnote{网的边定义为分支点与接点的补集的连通分支。}都满足式 (22.4) 中角为 $\vartheta$ 的局部 BPS 条件。

存在 $ij$-、$jk$- 与 $ik$-弦相遇的接点，其中 $i\ne j\ne k\in\{1,\dots,N\}$（对 $D$ 型或 $E$ 型的 S 类理论，还可能发生其他类型的接点）。一般的接点见图 22.8。

\Fig[0.45\linewidth]{figures/fig_22_8.png}{22.8}{一般的弦接点（$A$ 型）。}

分支点与分支割线都由无序整数对 $(ij)$ 标号。为使网是有限的，其所有边必须或者闭合，或者起止于分支点或接点。三个有限 BPS 网的例子见图 22.9。

% [原文 p.382]
任何这样的有限网都对应于 Seiberg–Witten 曲线 $\Sigma$ 上的一个圈 $C$，而 BPS 态的电荷就是 $C$ 的同调类：电荷格为 $\Gamma=H_1(\Sigma,\Z)$。$4d$ BPS 态的简并度编码在式 (22.7) 定义的第二螺旋度超迹 $\Omega(\gamma;u)$ 中，其中 $\gamma\in\Gamma$，$u$ 是库仑分支 $\mathcal{C}$ 上的点。

\Fig[0.95\linewidth]{figures/fig_22_9.png}{22.9}{三个有限 BPS 网（分支点用红色叉号表示）。}

\subsubsection*{面缺陷与孤子}

\begin{proposition}[22.16]
对应于 $A$、$D$ 或 $E$ 型李代数 $\mathfrak{g}$ 的 S 类理论容许典范的半 BPS 面缺陷 $\mathbb{S}_z$，它依赖于点 $z\in C_{g,n}$ 与 $\mathfrak{g}$ 的一个表示 [147]。
\end{proposition}

我们关注该表示是 $\mathfrak{su}(N)$ 基本表示的情形。在 $A_{N-1}$ 型 S 类理论的 M 理论实现中，这些面缺陷对应于开 M2-膜，其边界沿 $\R^{1,3}$ 延伸，并定域在 $C_{\mathrm{UV}}$ 中的点 $z$ 处，从而在低能四维 S 类理论中定义一个面缺陷。

在点 $u\in\mathcal{C}$ 处（它指定 Seiberg–Witten 覆叠 $\Sigma\to C_{\mathrm{UV}}$），$z$ 在 $\Sigma$ 中的原像对应于面缺陷 $\mathbb{S}_z$ 的不同真空：若 $z$ 是纤维化 $\Sigma\to C_{\mathrm{UV}}$ 的一般点，则面缺陷 $\mathbb{S}_z$ 有 $N$ 个不同的真空。在面缺陷 $\mathbb{S}_z$ 存在时，理论支持束缚在该面缺陷上、并在 $\mathbb{S}_z$ 的不同真空之间插值的 BPS 粒子。

\begin{definition}[22.17]
这些 BPS 粒子称为孤子，或 $2d$-$4d$ BPS 态。
\end{definition}

若 $z$ 是一般的，则在 $\mathbb{S}_z$ 的第 $i$ 个真空 $z^{(i)}$ 与第 $j$ 个真空 $z^{(j)}$（其中 $i\ne j$）之间插值的孤子称为 $ij$-孤子。孤子也可以描述为 $C_{\mathrm{UV}}$ 上终止于 $z\in C_{\mathrm{UV}}$ 的有限弦网。面缺陷 $\mathbb{S}_z$ 的三个 $12$-孤子的例子见图 22.10。

$C_{\mathrm{UV}}$ 上对应于 $\mathbb{S}_z$ 的 $ij$-孤子的弦网到 $\Sigma$ 的提升，是一条边界为 $z^{(j)}-z^{(i)}$ 的 1-链。记 $\Gamma_{ij}(z,z)$ 为 $\Sigma$ 上边界为 $z^{(j)}-z^{(i)}$ 的相对同调

% [原文 p.383]
类的集合，则 $\mathbb{S}_z$ 的孤子电荷属于集合
\[
\Gamma(z,z)=\bigsqcup_{i\ne j}\Gamma_{ij}(z,z).
\]
电荷为 $a\in\Gamma(z,z)$ 的孤子的中心荷 $Z_a(z)$ 可以通过沿任意代表 $a$ 的 1-链积分 Liouville 1-形式 $\lambda$ 来计算，与 $4d$ BPS 态的情形类似。此外，远离边缘稳定性墙时，存在一个受保护指标，计数电荷为 $a$ 的孤子（受制于一个「模 2 的注意事项」，稍后讨论）：
\begin{equation}
\mu(a)=\operatorname{Tr}_{\mathcal{H}^{1,\mathrm{BPS}}_{\mathbb{S}_z,a}(u)}F\ee^{\ii\pi F},\tag{22.11}
\end{equation}
其中 $F$ 是费米子数算子。

\Fig[0.95\linewidth]{figures/fig_22_10.png}{22.10}{面缺陷 $\mathbb{S}_z$ 的三个 $12$-孤子。}

\subsubsection*{带标架的 $2d$-$4d$ BPS 态}

在两个面缺陷 $\mathbb{S}_{z_1}$ 与 $\mathbb{S}_{z_2}$（其中 $z_1,z_2\in C_{\mathrm{UV}}$）存在时，可以如下定义额外的 BPS 态。可以考虑面缺陷 $\mathbb{S}_{z_1}$ 与 $\mathbb{S}_{z_2}$ 之间的超对称界面。若 $\R^{1,3}$ 按通常方式由 $x^{0,1,2,3}$ 坐标化，令 $\mathbb{S}_{z_1}$ 沿 $x^0$ 与 $x^1<0$ 延伸，而定域在 $x^{2,3}=0$；令 $\mathbb{S}_{z_2}$ 沿 $x^0$ 与 $x^1>0$ 延伸，也定域在 $x^{2,3}=0$。这些缺陷之间的界面沿 $x^0$ 延伸并定域在 $x^{1,2,3}=0$。

$\mathbb{S}_{z_1}$ 与 $\mathbb{S}_{z_2}$ 之间不同的界面由 $C_{\mathrm{UV}}$ 上连接 $z_1$ 与 $z_2$ 的路径 $\wp$ 的同伦类以及一个实参数 $\vartheta$ 分类。这些界面记作 $L_{\wp,\vartheta}$，它们是半 BPS 的，意即它们保持 $\mathbb{S}_{z_1}$ 与 $\mathbb{S}_{z_2}$ 所保持的四个超对称中的两个 [152]。

\begin{definition}[22.18]
对应于超对称界面 $L_{\wp,\vartheta}$ 的带标架 $2d$-$4d$ BPS 态，是保持复合系统的两个超荷的 Hilbert 空间 $\mathcal{H}^1_{L_{\wp,\vartheta}}$ 中的一个态。
\end{definition}

\begin{remark}[22.19]
若不存在 $\gamma\in\Gamma$ 使得 $\ee^{-\ii\vartheta}Z_\gamma\in\R_-$，则称参数 $\vartheta$ 是一般的。类似地，若不存在 $a\in\Gamma(z_1,z_2)\cup$

% [原文 p.384]
$\Gamma(z_2,z_2)$ 使得 $\ee^{-\ii\vartheta}Z_a\in\R_-$，则称 $(\wp,\vartheta)$ 是一般的。当 $(\wp,\vartheta)$ 一般时，对应于 $(\wp,\vartheta)$ 的带标架 $2d$-$4d$ BPS 态定域在 $L_{\wp,\vartheta}$ 附近。
\end{remark}

带标架 $2d$-$4d$ BPS 态的电荷由边界为 $z_2^{(j)}-z_1^{(i)}$ 的 1-链的相对同调类分类，其中 $z_2^{(j)}$（相应地 $z_1^{(i)}$）是 $z_2$（相应地 $z_1$）的第 $j$ 个（相应地第 $i$ 个）提升。记 $\Gamma_{ij}(z_1,z_2)$ 为这类同调类的集合，并定义
\[
\Gamma(z_1,z_2)=\bigsqcup_{i,j}\Gamma_{ij}(z_1,z_2).
\]
还存在另一个受保护指标，称为带标架 BPS 指标，它计数带标架 $2d$-$4d$ BPS 态 [152]：
\begin{equation}
\overline{\Omega}{}'(L_{\wp,\vartheta},a)=\operatorname{Tr}_{\mathcal{H}^{1,\mathrm{BPS}}_{L_{\wp,\vartheta},a}}\ee^{\ii\pi F},\tag{22.12}
\end{equation}
其中 $a\in\Gamma(z_1,z_2)$。

\subsubsection*{$\Z_2$-扭转}

式 (22.12) 中的算子 $F$ 只定义到相差复数的平移。要求 $\overline{\Omega}{}'(L_{\wp,\vartheta},a)$ 为实，这限制了 $F$，但仍允许整数的平移。因此，$\overline{\Omega}{}'(L_{\wp,\vartheta},a)$ 仍然只定义到相差一个符号。类似地，式 (22.11) 中的指标依赖于费米子数算子 $F$ 的选取，这也使它只确定到相差一个符号。

如 [154] 所提出，有一种系统的方法可以在整体上消解这些含混性，它类似于 [129] 中定义扭曲局部系统的模空间 $\mathcal{M}_{GL,S}$ 时所用的扭转（另见注 10.21）。对任意光滑曲面 $S$，记 $\widetilde{S}$ 为其单位切丛，$H_S$ 为圆丛 $\widetilde{S}\to S$ 的纤维类。注意 $S$ 上任意光滑有向路径都典范地提升到 $\widetilde{S}$，而 $S$ 上光滑无向路径在 $\widetilde{S}$ 中有两个典范提升。

现在考虑 $\Sigma$ 上边界为 $z^{(j)}-z^{(i)}$ 的 1-链。回忆 $\Gamma_{ij}(z,z)$ 表示这类 1-链的同调类的集合，对应于缺陷 $\mathbb{S}_z$ 的 $ij$-孤子的电荷。设 $z_+^{(i)},z_+^{(j)}\in\widetilde{\Sigma}$ 表示 $z^{(i)}$ 与 $z^{(j)}$ 到 $\widetilde{\Sigma}$ 的提升，其定向与该 1-链的定向一致。通过考虑集合 $\widetilde{\Gamma}(z_+^{(i)},z_+^{(j)})$ 中 $ij$-孤子的电荷——该集合由 $\widetilde{\Sigma}$ 中边界为 $z_+^{(j)}-z_+^{(i)}$ 的相对同调类模 $2H_\Sigma$ 组成——而不是考虑 $\Gamma_{ij}(z,z)$，可以使孤子简并度 $\mu$ 变得完全良定义。这通过如下条件实现：对 $\overline{a}\in\Gamma_{ij}(z,z)$ 在 $\widetilde{\Gamma}(z_+^{(i)},z_+^{(j)})$ 中的两个提升 $a$ 与 $a'$，有关系
\[
\mu(a)/\mu(a')=(-1)^{w(a,a')},
\]
其中 $w(a,a')$ 是 $a$ 与 $a'$ 在绕数上的差（以 $H_\Sigma$ 为单位，模 $2H_\Sigma$）。令

% [原文 p.385]
\[
\widetilde{\Gamma}(\widetilde{z},\widetilde{z})=\bigsqcup_{i\ne j}\widetilde{\Gamma}(z_+^{(i)},z_+^{(j)}).
\]

\begin{definition}[22.20]
对 $a\in\widetilde{\Gamma}(z_+^{(i)},z_+^{(j)})$，整数 $\mu(a)$ 称为孤子简并度。
\end{definition}

类似地，带标架 BPS 指标也可以通过考虑 $\overline{a}$ 到 $\widetilde{\Sigma}$ 的提升 $a$（模 $2H_\Sigma$）而变得良定义。确切地说，代替类 $\overline{a}\in\Gamma(z_1,z_2)$，考虑 $\overline{a}$ 到 $\widetilde{\Sigma}$ 的提升 $a$，即边界为 $z_{2,+}^{(j)}-z_{1,+}^{(i)}$（对某对 $i,j$）的相对同调类，其中 $+$ 表示由 $\wp$ 规定的 $z_2^{(j)}$ 与 $z_1^{(i)}$ 到 $\widetilde{\Sigma}$ 的提升，模 $2H_\Sigma$。它们构成集合 $\widetilde{\Gamma}(\widetilde{z}_1,\widetilde{z}_2)$。两个这样的提升 $a$ 与 $a'$ 的相对同调相差一个绕数 $w(a,a')$，它取 $0$ 或 $1$。我们规定
\[
\overline{\Omega}{}'(L_{\wp,\vartheta},a)/\overline{\Omega}{}'(L_{\wp,\vartheta},a')=(-1)^{w(a,a')}.
\]
同样地，给定两条路径 $\wp$ 与 $\wp'$，其到 $\widehat{C}_{\mathrm{UV}}$ 的提升在 $\widehat{C}_{\mathrm{UV}}$ 中具有相同的起点与终点，则可以定义绕数 $w(\wp,\wp')$，如图 22.11 所示，并规定
\[
\overline{\Omega}{}'(L_{\wp,\vartheta},a)/\overline{\Omega}{}'(L_{\wp',\vartheta},a)=(-1)^{w(\wp,\wp')}.
\]
这意味着界面 $L_{\wp,\vartheta}$ 的带标架 BPS 指标 $\overline{\Omega}{}'$ 是路径 $\wp$ 的扭化同伦不变量，而非真正的同伦不变量。

\Fig[1.0\linewidth]{figures/fig_22_11.png}{22.11}{两条路径 $\wp$ 与 $\wp'$，使得 $w(\wp,\wp')=1$。}

% [Tutor Note: $\Z_2$-扭转的要点可这样核对：(22.11) 与 (22.12) 里的 $F$ 只定义到「整数平移」，而把 $F$ 平移 $c\in\Z$ 会使 $\ee^{\ii\pi F}$ 变号 $(-1)^c$；另一方面绕数只在模 $2H_\Sigma$ 的意义下有意义，$w(a,a')\in\{0,1\}$。两者恰好互相抵消：把 $a$ 换成绕数加一的另一提升时，指标变号，而 $F$ 的整数平移也恰好给出同一变号，故定义 $\mu(a)/\mu(a')=(-1)^{w(a,a')}$ 后 $\mu$ 在整个上同调类上是一致的。若把 $2H_\Sigma$ 换成 $H_\Sigma$（即模 $1$ 而非模 $2$），则 $(-1)^w$ 不再是良定义的配对，这正是必须用 $2H_\Sigma$ 而非 $H_\Sigma$ 的原因。]
\subsubsection*{带标架 BPS 态的生成函数}
给定 $C_{\mathrm{UV}}$ 中两点 $z_1$ 与 $z_2$ 之间的路径 $\wp$，带标架 BPS 指标 $\overline{\Omega}{}'(L_{\wp,\vartheta},a)$ 可以如下编码为一个形式生成函数 $F(\wp,\vartheta)$。对每个 $a\in\widetilde{\Gamma}(\widetilde{z}_1,\widetilde{z}_2)$，令 $X_a$ 表示一个形式变量，并满足条件：若 $a$ 与 $a'$ 投影到同一个类 $\overline{a}\in\Gamma(z_1,z_2)$，则
\[
X_a/X_{a'}=(-1)^{w(a,a')}.
\]
这允许我们定义
\begin{equation}
F(\wp,\vartheta)=\sum_{\overline{a}\in\Gamma(z_1,z_2)}\overline{\Omega}{}'(L_{\wp,\vartheta},a)X_a,\tag{22.13}
\end{equation}

% [原文 p.386]
其中 $a$ 是 $\overline{a}$ 到 $\widetilde{\Gamma}(\widetilde{z}_1,\widetilde{z}_2)$ 的任一提升，因为对每个 $\overline{a}$，乘积 $\overline{\Omega}{}'(L_{\wp,\vartheta},a)X_a$ 与代表 $a$ 的选取无关。$\widetilde{\Gamma}(\widetilde{z}_1,\widetilde{z}_2)$ 的性质蕴含 $F(\wp,\vartheta)$ 是路径 $\wp$ 的（扭化）同伦不变量。

如 13.2 节所讨论，只要第一条路径的终点与第二条路径的起点重合，就可以把 $\widetilde{\Sigma}$ 上的路径连接起来。这给出 $\widetilde{\Sigma}$ 上的（同调）路径代数，其乘法规则如下：
\[
X_aX_b=\begin{cases}
X_{a+b} & \text{若 } \operatorname{end}(a)=\operatorname{start}(b),\\
0 & \text{否则}.
\end{cases}
\]

\begin{proposition}[22.21]
若 $\wp$ 与 $\wp'$ 是 $C_{\mathrm{UV}}$ 上两条路径，且 $\operatorname{end}(\wp)=\operatorname{start}(\wp')$（在 $\widehat{C}_{\mathrm{UV}}$ 中），即 $\wp$ 与 $\wp'$ 可以连接成一条光滑路径 $\wp\wp'$，并且 $(\wp,\vartheta)$ 与 $(\wp',\vartheta)$ 都是一般的，则
\[
F(\wp,\vartheta)F(\wp',\vartheta)=F(\wp\wp',\vartheta).
\]
\end{proposition}

\begin{proof}
这一结果来自对 $(\wp,\vartheta)$ 与 $(\wp',\vartheta)$ 的一般性假设。在该假设下，带标架 $2d$-$4d$ BPS 态分别定域在界面 $L_{\wp,\vartheta}$ 与 $L_{\wp',\vartheta}$ 附近。因此，乘积 $F(\wp,\vartheta)F(\wp',\vartheta)$ 不依赖于 $L_{\wp,\vartheta}$ 与 $L_{\wp',\vartheta}$ 之间的距离。在此距离趋于零的极限情形中，这些带标架 $2d$-$4d$ BPS 态正是界面 $L_{\wp\wp',\vartheta}$ 的那些态。
\end{proof}

我们关心的是计算生成函数 $F(\wp,\vartheta)$。这一计算关键地依赖于 $\wp$ 与 $C_{\mathrm{UV}}$ 上的谱网络 $\mathcal{W}_\vartheta$ 相交的方式，下一节将详细说明。

\subsection{带标架 $2d$-$4d$ BPS 指标的计算}

\begin{definition}[22.22]
设 $\vartheta$ 是一个角。若 $Z_a(z)/\ee^{\ii\vartheta}\in\R_-$，则称点 $z\in C_{\mathrm{UV}}$ 支持孤子电荷 $a\in\Gamma(z,z)$。谱网络 $\mathcal{W}_\vartheta$ 是支持某个满足 $\mu(a)\ne0$ 的 $a\in\Gamma(z,z)$ 的点 $z\in C_{\mathrm{UV}}$ 的集合，其中 $a$ 是 $\overline{a}$ 到 $\widetilde{\Gamma}(\widetilde{z},\widetilde{z})$ 的任一提升。
\end{definition}

\begin{proposition}[22.23]
谱网络 $\mathcal{W}_\vartheta$ 在 $C_{\mathrm{UV}}$ 中余维为 1，由若干线段之并组成，这些线段称为 S-壁，终止于分支点或称为接点的特殊点，不同的 S-壁在接点处相交。
\end{proposition}

我们假定当 $\vartheta$ 一般时，$\mathcal{W}_\vartheta$ 的任一一般点恰支持某个 $i\ne j$ 的 $a\in\Gamma_{ij}(z,z)$，其中 $\mathcal{W}_\vartheta$ 的一般点指 $\mathcal{W}_\vartheta$ 中既不是分支点也不是接点的所有点。此时，给定 S-壁 $p$ 的所有点都支持同一个电荷 $a_{ij}\in\Gamma(z,z)$。称该壁为 $ij$ 型，或说它是一条 $ij$-轨道，因为按定义这意味着对所有 $z\in p$ 都有 $Z_a/\ee^{\ii\vartheta}\in\R_-$，进而 $p$ 满足式 (22.4) 中角为 $\vartheta$ 的局部 BPS 微分方程，这是因为

% [原文 p.387]
\[
\dd Z_a(z)=\frac{1}{\pi}\big(\lambda^{(j)}-\lambda^{(i)}\big).
\]
换言之，S-壁满足与 22.1.2 节所述 $4d$ BPS 弦相同的微分方程。

\subsubsection*{质量滤过}

给定具有质量量纲的参数 $\Lambda$，可以定义截断 $\mathcal{W}_\vartheta[\Lambda]$ 为 $\mathcal{W}_\vartheta$ 中仅含那些满足「$z$ 支持某个使 $|Z_a(z)|<\Lambda$ 的电荷 $a\in\Gamma(z,z)$」的点 $z\in\mathcal{W}_\vartheta$ 的子集。对任意两个质量参数 $\Lambda<\Lambda'$，包含关系 $\mathcal{W}_\vartheta[\Lambda]\subset\mathcal{W}_\vartheta[\Lambda']$ 成立，这为「质量滤过」这一名称提供了依据。类似地，可以如下定义截断的孤子简并度：
\[
\mu[\Lambda]:\ \big\{a\in\widetilde{\Gamma}(\widetilde{z},\widetilde{z})\ \big|\ |Z_a(z)|<\Lambda\big\}\longrightarrow \mu[\Lambda](a):=\mu(a).
\]
对 $\Lambda=0$，截断网络 $\mathcal{W}_\vartheta[\Lambda]$ 只包含覆叠 $\Sigma\to C_{\mathrm{UV}}$ 的分支点。当 $\Lambda$ 从 $0$ 逐渐增大到 $\infty$ 时，网络逐步生长，最终在 $\Lambda\to\infty$ 的极限中重构出整个谱网络：
\[
\mathcal{W}_\vartheta=\lim_{\Lambda\to\infty}\mathcal{W}_\vartheta[\Lambda].
\]
对充分小的 $\Lambda$，截断网络 $\mathcal{W}_\vartheta[\Lambda]$ 只由从覆叠 $\Sigma\to C_{\mathrm{UV}}$ 的分支点发出的小 S-壁组成（参见 22.1.2 节）。当 $\Lambda$ 增大时，某些 S-壁可能相交，形成接点。当 $\vartheta$ 一般时，这样的接点属于以下类型之一：
\begin{itemize}
\item 当 $ij$ 型 S-壁与另一条 $kl$ 型 S-壁相交（其中指标 $i,j,k,l$ 两两不同）时，相互作用是平凡的：两条 S-壁只是彼此穿过，并沿各自的轨道继续延伸。
\item 当 $ij$ 型 S-壁与另一条 $jk$ 型 S-壁相交（其中 $k\ne i$）时，接点处可能产生一条 $ik$ 型 S-壁。
\item 反之，$ik$ 型 S-壁可以在它与另外两条分别属于 $ij$ 型与 $jk$ 型的 S-壁相遇的接点处消失。
\item 也可能存在 $ij$、$jk$ 与 $ik$ 型 S-壁相遇并交叉的接点（参见图 22.8）。
\end{itemize}

\begin{remark}[22.24]
壁在接点处产生或消失的可能性与谱网络的定义 22.22 相关。按该定义，S-壁并不由支持给定孤子电荷的全部潜在轨道组成，而只由那些孤子简并度 $\mu(a)$ 非零的轨道组成。这一区别将在随后几节中进一步探讨。
\end{remark}

% [原文 p.388]
\subsubsection*{$F(\wp,\vartheta)$ 的计算}

$F(\wp,\vartheta)$ 的计算来自 [152] 中导出的 $2d$-$4d$ 穿墙公式。给定路径 $\wp\subset C_{\mathrm{UV}}$，先定义 $\Sigma$ 上同调路径代数的元素
\begin{equation}
\mathcal{D}(\wp):=\sum_{i=1}^{N}X_{\wp^{(i)}}\tag{22.14}
\end{equation}
其中 $\wp^{(i)}$，$i=1,\dots,N$，是 $\wp$ 到 $\Sigma$ 的提升。

由命题 22.21，在假定 $\wp$ 只在一般点处与谱网络 $\mathcal{W}_\vartheta$ 相交的前提下，$F(\wp,\vartheta)$ 的计算完全由两种情形决定：
\begin{enumerate}
\item 对不与 $\mathcal{W}_\vartheta$ 相交的小路径 $\wp$：
\begin{equation}
F(\wp,\vartheta)=\mathcal{D}(\wp),\tag{22.15}
\end{equation}
其中 $\mathcal{D}(\wp)$ 由式 (22.14) 定义。
\item 对在一般点处恰好与 $\mathcal{W}_\vartheta$ 相交一次的小路径 $\wp$，假定 $\wp$ 在与 $\mathcal{W}_\vartheta$ 的交点 $z$ 处落在一条携带孤子电荷 $\overline{a}\in\Gamma_{ij}(z,z)$ 的 S-壁上。记孤子重数为 $\mu(a)$，其中 $a$ 是 $\overline{a}$ 到 $\widetilde{\Gamma}(\widetilde{z},\widetilde{z})$ 的任一提升。在这种情况下，如图 22.12 所示把 $\wp$ 分成两条路径 $\wp_+$ 与 $\wp_-$。
\end{enumerate}

\Fig[0.9\linewidth]{figures/fig_22_12.png}{22.12}{当 $\wp$ 与 $\mathcal{W}_\vartheta$ 相交一次时 $F(\wp,\vartheta)$ 的计算。}

于是生成函数 $F(\wp,\vartheta)$ 表为
\[
F(\wp,\vartheta)=\mathcal{D}(\wp)+\mathcal{D}(\wp_+)\big(\mu(a)X_{a_{\widetilde{z}}}\big)\mathcal{D}(\wp_-),
\]
其中 $a_{\widetilde{z}}$ 是 $\widetilde{\Sigma}$ 中的路径：它由从 $\widetilde{z}_+^{(i)}$ 到 $a$ 的起点的最短路径、路径 $a$ 本身、以及从 $a$ 的终点到 $\widetilde{z}_-^{(j)}$ 的最短路径连接而成。这里 $\widetilde{z}_+^{(i)}$（相应地 $\widetilde{z}_-^{(j)}$）表示 $z^{(i)}$（相应地 $z^{(j)}$）到 $\widetilde{\Sigma}$ 的、由 $\wp_+$（相应地 $\wp_-$）规定的提升。注意该式第二项中唯一的非零贡献是
\[
\mu(a)X_{\wp^{(i)}_+}X_{a_{\widetilde{z}}}X_{\wp^{(j)}_-},
\]
因为其他路径无法连接。于是，穿过一条 S-壁给 $F(\wp,\vartheta)$ 引入一项额外贡献，它对应于 $\wp$ 沿与壁上所支持 BPS 孤子相伴的路径的一次绕行。

最后，$F(\wp,\vartheta)$ 可以改写为

% [原文 p.389]
\begin{equation}
F(\wp,\vartheta)=\mathcal{D}(\wp_+)\big(1+\mu(a)X_{a_{\widetilde{z}}}\big)\mathcal{D}(\wp_-).\tag{22.16}
\end{equation}

\subsection{同伦不变性与孤子简并度}

对带标架 BPS 指标 $\overline{\Omega}{}'(L_{\wp,\vartheta},a)$、从而对生成函数 $F(\wp,\vartheta)$ 的一个重要约束是：它们应当是路径 $\wp$ 的（扭化）同伦不变量。这又使得可以确定谱网络 $\mathcal{W}_\vartheta$ 中每条 S-壁的孤子简并度 $\mu(a)$。本节中我们再次假定 $\vartheta$ 是一般的。

\begin{proposition}[22.25]
对 $\overline{a}$ 的任一提升 $a$，孤子简并度 $\mu(a)$ 沿每条 S-壁为常数。因此，若 $p$ 是一条携带电荷 $\overline{a}\in\Gamma(z,z)$ 的 S-壁，则记 $\mu(a,p)$ 为对应于 $\overline{a}$ 的提升 $a$ 沿 $p$ 的孤子简并度。
\end{proposition}

\begin{proof}
考虑图 22.13 的两个设定。由于 $\wp'=\wp'_+\wp'_0\wp'_-$ 与 $\wp$ 属于同一个扭化同伦类，必有
\[
F(\wp,\vartheta)=F(\wp',\vartheta).
\]
由式 (22.15) 有

\Fig[0.75\linewidth]{figures/fig_22_13.png}{22.13}{两个（正则的）同伦等价设定。}

\[
F(\wp,\vartheta)=\mathcal{D}(\wp),
\]
而式 (22.16) 蕴含 $F(\wp',\vartheta)$ 分解为
\begin{align*}
&\mathcal{D}(\wp'_+)\big(1+\mu(a;z_+)X_{a_{\widetilde{z}_+}}\big)\mathcal{D}(\wp'_0)\big(1+\mu(a;z_-)X_{a_{\widetilde{z}_-}}\big)\mathcal{D}(\wp'_-)\\
={}&\mathcal{D}(\wp')+\mu(a;z_+)X_{(\wp'_+)^{(i)}a_{\widetilde{z}_+}(\wp'_0)^{(j)}(\wp'_-)^{(j)}}+\mu(a;z_-)X_{(\wp'_+)^{(i)}(\wp'_0)^{(i)}a_{\widetilde{z}_-}(\wp'_-)^{(j)}}\\
={}&\mathcal{D}(\wp')+\big(\mu(a;z_+)-\mu(a;z_-)\big)X_{(\wp'_+)^{(i)}a_{\widetilde{z}_+}(\wp'_0)^{(j)}(\wp'_-)^{(j)}},
\end{align*}
其中最后一个等号来自如下事实：路径 $(\wp'_+)^{(i)}a_{\widetilde{z}_+}(\wp'_0)^{(j)}(\wp'_-)^{(j)}$ 与 $(\wp'_+)^{(i)}(\wp'_0)^{(i)}a_{\widetilde{z}_-}(\wp'_-)^{(j)}$ 只相差绕 $\widetilde{\Sigma}\to\Sigma$ 的圆纤维一周的绕数单位。因此，同伦不变性要求 $\mu(a;z_+)=\mu(a;z_-)$，从而 $\mu(a)$ 沿该 S-壁为常数。
\end{proof}

\begin{proposition}[22.26]
从分支点发出的任意 S-壁的孤子简并度为 $1$。
\end{proposition}

% [原文 p.390]
\begin{proof}
这同样来自图 22.14 所示设定中的同伦不变性，图中用紫色给出了一个分支割线的选取。生成函数

\Fig[0.42\linewidth]{figures/fig_22_14.png}{22.14}{分支点附近的同伦不变性。}

$F(\wp,\vartheta)$ 可以直接由式 (22.16) 表出，只是分支割线把 $\wp^{(i)}_0$ 与 $\wp^{(j)}_-$ 相连、反之亦然。另一方面，生成函数 $F(\wp',\vartheta)$ 如前一命题的证明那样应用式 (22.16) 两次而导出。

注意在分支点附近，每条 S-壁上所支持的孤子由该 S-壁上连接所考虑的点与分支点的那一段表示。这意味着对每条 S-壁都有 $\mu(a)=\pm1$。直接验证表明，同伦不变性强加 $\mu(a)=1$。
\end{proof}

\Fig[0.5\linewidth]{figures/fig_22_15.png}{22.15}{一个一般接点。}

\begin{proposition}[22.27]
考虑图 22.15 所示的一般接点。则
\begin{align*}
\mu(a,p')&=\mu(a,p),\\
\mu(b,r')&=\mu(b,r),\\
\mu(c,q')&=\mu(c,q)+(-1)^{w(a_{\widetilde{z}}+b_{\widetilde{z}},c_{\widetilde{z}})}\mu(a,p)\mu(b,r).
\end{align*}
\end{proposition}

\begin{proof}
这同样来自同伦不变性。为证明该命题，把式 (22.16) 应用三次：先对一条从 $z$ 左侧竖直通过的路径 $\wp$，然后对另一条从 $z$ 右侧竖直通过的路径 $\wp'$。这些路径随后可以形变，使它们与 S-壁的交点只出现在接点 $z$ 处。
\end{proof}

命题 22.27 一般公式的重要特例是 $\mu(c,q)=0$ 而 $\mu(c,q')\ne0$ 的情形（对应于 $ik$-壁在

% [原文 p.391]
$ij$-$jk$ 接点处产生），以及 $\mu(c,q')=0$ 而 $\mu(c,q)\ne0$ 的情形（对应于 $ik$-壁在 $ij$-$jk$-$ik$ 接点处湮灭）。

\subsection{由谱网络得到 $4d$ BPS 简并度}

BPS 谱由考察带标架 $2d$-$4d$ BPS 态的生成函数 $F(\wp,\vartheta)$ 随 $\vartheta$ 变化的方式而确定。一般地，$F(\wp,\vartheta)$ 在 $\vartheta$ 的小变化下保持常数。然而，存在某些特定的 $\vartheta$ 值，使 $F(\wp,\vartheta)$ 发生不连续变化。这可能出于两个不同的原因。

第一，当 $\vartheta$ 变化时，谱网络 $\mathcal{W}_\vartheta$ 的 S-壁在 $C_{\mathrm{UV}}$ 上移动。由于这种移动，S-壁可能穿过所考虑路径 $\wp$ 的起点或终点，从而按式 (22.16) 所述增加或去掉一项贡献。这给出 $F(\wp,\vartheta)$ 跳跃的一个原因，它既依赖于 $\vartheta$ 也依赖于 $\wp$ 的端点。

$F(\wp,\vartheta)$ 在 $\vartheta$ 的临界值 $\vartheta_c$ 处不连续的第二个原因在于 $\mathcal{W}_\vartheta$ 拓扑的改变。如 22.2.4 节所讨论，这类改变发生在 $\vartheta$ 的特定值处，这些值称为 $\mathcal{K}$-壁，此时一条 $ij$-轨道遇到一个 $(ij)$-分支点。这种拓扑变化最简单的例子称为鞍点连接，见图 22.16。

\Fig[0.98\linewidth]{figures/fig_22_16.png}{22.16}{谱网络 $\mathcal{W}_\vartheta$ 在鞍点连接处的拓扑变化。}

图 22.17 给出了这类连接的一个实例，它显示纯 $SU(2)$ SYM 理论在库仑分支原点处、对不同的 $\vartheta$ 值的谱网络。在此例中，$\mathcal{K}$-壁出现在 $\theta_c=\pi\sim0$ 与 $\theta_c=\pi/2$ 处。

\subsubsection*{$\mathcal{K}$-壁公式}

一般地，在 $\vartheta=\vartheta_c$ 处，谱网络 $\mathcal{W}_{\vartheta_c}$ 表现出一条非一般的 S-壁，它由两条或多条一般 S-壁的碰撞形成。当这些 S-壁沿相反方向行进时（鞍点连接就是这种情形），所得的非一般 S-壁称为双向街或双壁。

% [原文 p.392]
\Fig[0.75\linewidth]{figures/fig_22_17.png}{22.17}{纯 $SU(2)$ SYM 理论在库仑分支原点处的谱网络随 $\vartheta$ 从 $0$ 变到 $\pi$ 的演化。}

每条双向街都容许两种不同的消解，取决于 $\vartheta=\vartheta_c+\epsilon$ 还是 $\vartheta=\vartheta_c-\epsilon$，其中 $\epsilon>0$ 无穷小。若 $\wp$ 是一条穿过双向街的路径，则可以用式 (22.16) 计算 $F(\wp,\vartheta_c+\epsilon)$ 与 $F(\wp,\vartheta_c-\epsilon)$，结果发现 $F(\wp,\vartheta_c+\epsilon)\ne F(\wp,\vartheta_c-\epsilon)$：生成函数 $F(\wp,\vartheta)$ 在 $\vartheta=\vartheta_c$ 处发生跳跃。

生成函数 $F(\wp,\vartheta)$ 在 $\vartheta$ 穿过 $\mathcal{K}$-壁时的跳跃，可以表为作用在形式路径变量 $X_a$ 上的一个普适替换 $\mathcal{K}$：
\begin{equation}
F(\wp,\vartheta_c+\epsilon)=\mathcal{K}\big[F(\wp,\vartheta_c-\epsilon)\big].\tag{22.17}
\end{equation}
式 (22.17) 称为 $\mathcal{K}$-壁公式。

双向街 $p$ 上的孤子简并度汇集为一个生成函数 $Q(p)$，为此必须为 $\widetilde{\Sigma}$ 上闭路径的每个同调类 $\widetilde{\gamma}$（即 $\widetilde{\gamma}\in\widetilde{\Gamma}=H_1(\widetilde{\Sigma},\Z)$）模 $2H_\Sigma$ 的平移引入新变量 $X_{\widetilde{\gamma}}$。同调路径代数随后扩充到包含生成元 $X_{\widetilde{\gamma}}$，其乘法规则为
\[
X_{\widetilde{\gamma}}X_{\widetilde{\gamma}'}=X_{\widetilde{\gamma}+\widetilde{\gamma}'}\,,\qquad X_{\widetilde{\gamma}}X_a=X_{\widetilde{\gamma}+a},
\]
其中 $\widetilde{\gamma}+a$ 表示 $\widetilde{\Gamma}$ 对 $\widetilde{\Gamma}(\widetilde{z},\widetilde{z})$ 的自然作用（对任意 $z\in C_{\mathrm{UV}}$）。此外还有附加关系 $X_{H_\Sigma}=-1$。

此外，对任一开路径 $a\in\widetilde{\Gamma}_{ij}(\widetilde{z},\widetilde{z})$，记 $\mathrm{cl}(a)\in\widetilde{\Gamma}$ 为相应的闭同调类。用此记号，生成函数 $Q(p)$

% [原文 p.393]
定义为
\[
Q(p)=1+\sum_{\overline{a}\in\Gamma_{ij}(p),\ \overline{b}\in\Gamma_{ji}(p)}\mu^-(a,p)\mu^-(b,p)X_{\mathrm{cl}(a+b)},
\]
其中 $\Gamma_{ij}(p)$ 与 $\Gamma_{ji}(p)$ 表示双向街 $p$ 所支持的两组电荷，而 $\mu^-(a,p)$ 与 $\mu^-(b,p)$ 表示在 $\vartheta\to\vartheta_c^-$ 极限中 $p$ 上的孤子简并度。把 $\mu^-$ 替换为 $\mu^+$（其中 $\mu^+$ 在 $\vartheta\to\vartheta_c^+$ 极限中定义）得到同一个生成函数 $Q(p)$。

注意 $Q(p)$ 中出现的所有 $X_{\widetilde{\gamma}}$ 都满足 $\ee^{\ii\vartheta}Z_{\widetilde{\gamma}}\in\R_-$。一般地，且典型地当考虑库仑分支上远离边缘稳定性墙的点时，可以假定使 $\ee^{\ii\vartheta}Z_\gamma\in\R_-$ 的电荷 $\gamma\in\Gamma$ 的集合 $\Gamma_c$ 由单个元素 $\gamma_0$ 生成。此时 $Q(p)$ 成为单变量 $X_{\widetilde{\gamma}_0}$ 的幂级数，其中 $\widetilde{\gamma}_0$ 是 $\gamma_0$ 到 $\widetilde{\Gamma}$ 的一个提升。事实上，对每个 $\gamma\in\Gamma$ 都存在一个「优选的」提升 $\widetilde{\gamma}\in\widetilde{\Gamma}$，使得 $Q(p)$ 分解为
\[
Q(p)=\prod_{\gamma\in\Gamma_c}\big(1-X_{\widetilde{\gamma}}\big)^{\alpha_\gamma(p)},
\]
其指数满足 $\alpha_\gamma(p)\in\Z$。于是，对每个 $\gamma\in\Gamma_c$ 定义
\[
L(\gamma)=\sum_{p}\alpha_\gamma(p)p_\Sigma,
\]
其中求和遍历 $\mathcal{W}_{\vartheta_c}$ 中的双向街，$p_\Sigma$ 表示 $p$ 在 $\Sigma$ 中的 1-链提升。可以证明 $L(\gamma)$ 事实上是一个 1-圈。这又使得可以定义 $L(\gamma)$ 与任意 $\overline{c}\in\Gamma(z_1,z_2)$ 的相交数 $\langle\overline{c},L(\gamma)\rangle$。由此得到如下结果 [154]：

\begin{proposition}[22.28]
普适替换 $\mathcal{K}$ 由下式给出：
\[
\mathcal{K}:\ X_a\longmapsto \prod_{\gamma\in\Gamma_c}\big(1-X_{\widetilde{\gamma}}\big)^{\langle a,L(\gamma)\rangle}X_a.
\]
\end{proposition}

% [Tutor Note: 可核对该公式与 (22.8) 的一致性。对鞍点连接（例 22.31），$\Gamma_c$ 由单个 $\gamma$ 生成，$L(\gamma)=p_\Sigma$ 且 $\Omega(\gamma)=1$，于是 $\mathcal{K}$ 化为 $X_a\mapsto X_a(1-X_{\widetilde{\gamma}})^{\langle a,p_\Sigma\rangle}$，与式 (22.8) 的 $K_\gamma$ 在取 $\sigma(\gamma)=1$、$\Omega(\gamma)=1$ 时的作用 $X_{\gamma'}\mapsto X_{\gamma'}(1-X_\gamma)^{\langle\gamma',\gamma\rangle}$ 完全同形。另可核对：例 22.32 的闭合圈给出 $L(\gamma)=-p^1_\Sigma-p^2_\Sigma$、$\Omega(\gamma)=-2$，此时 $\langle a,L(\gamma)\rangle$ 的符号与指数 $-2$ 恰好相互抵消，故该圈对 $\mathcal{K}$ 的作用等价于相应矢量多重态的贡献（系数 $-2$），与 22.3.4 节「闭合圈对应 BPS 矢量多重态」的论断相容。]

\subsubsection*{BPS 简并度}

我们注意到
\[
\ee^{-\ii\vartheta_c}\int_{p_\Sigma}\lambda\in\R_-.
\]
在 $\Gamma_c$ 由单个元素生成的假定下，这意味着 $L(\gamma)$ 的同调类 $[L(\gamma)]$ 与 $\gamma$ 成比例。

\begin{proposition}[22.29]
比值 $[L(\gamma)]/\gamma$ 就是第二螺旋度超迹 $\Omega(\gamma)$。
\end{proposition}

\begin{remark}[22.30]
事实上，谱网络不仅提供了计算 $4d$ BPS 简并度 $\Omega(\gamma)$ 的方法，还提供了计算增强 BPS 简并度 $\omega(\gamma,\overline{a})$ 的方法（对 $\gamma\in\Gamma$

% [原文 p.394]
与 $\overline{a}\in\Gamma(z_1,z_2)$）[152]，方式如下：
\[
\omega(\gamma,\overline{a})=\langle L(\gamma),\overline{a}\rangle.
\]
增强 BPS 简并度 $\omega(\gamma,\overline{a})$ 特别满足
\[
\omega(\gamma,\overline{a}+\gamma')-\omega(\gamma,\overline{a})=\Omega(\gamma)\langle\gamma,\gamma'\rangle,
\]
即它们编码了 $4d$ BPS 简并度 $\Omega(\gamma)$ 的信息。此外，与 $\Omega(\gamma)$ 不同，增强 BPS 简并度 $\omega(\gamma,\overline{a})$ 提供了关于 $4d$ BPS 态沿 $C_{\mathrm{UV}}$ 的位置及其与面缺陷相互作用的细节。
\end{remark}

\begin{example}[22.31]
在鞍点连接的情形中，构成双向街 $p$ 的两条 S-壁携带电荷 $\overline{a},\overline{b}$，使得 $\mu^-(a)=\mu^-(b)=1$，其中 $a$ 与 $b$ 是 $\overline{a}$ 与 $\overline{b}$ 的自然提升。记 $\widetilde{\gamma}$ 为 $\gamma=\mathrm{cl}(\overline{a}+\overline{b})$ 的优选提升，则有
\[
Q(p)=1-X_{\widetilde{\gamma}}.
\]
这意味着 $L(\gamma)=p_\Sigma$，其中 $[L(\gamma)]=\gamma$，且对任意 $n>1$ 有 $L(n\gamma)=0$。于是
\[
\Omega(\gamma)=1,
\]
这表明鞍点连接对应于 $4d$ BPS 超多重态。
\end{example}

\begin{example}[22.32]
$\mathcal{K}$-壁的另一个基本例子出现在绕圆柱缠绕的 S-壁塌缩为两条双向街 $p^1$ 与 $p^2$ 时。这两条双向街各自起止于同一个分支点。这一情形在 [154, 第 7.2 节] 中有描述。我们得到
\[
L(\gamma)=-p^1_\Sigma-p^2_\Sigma,
\]
其中 $p^1_\Sigma$ 与 $p^2_\Sigma$ 都在同调类 $[\gamma]$ 中。因此
\[
\Omega(\gamma)=-2,
\]
这表明谱网络中这样的闭合圈对应于 BPS 矢量多重态。
\end{example}

以类似方式可以分析在 $\mathcal{K}$-壁处出现的更复杂的有限弦网。特别地，可以计算相应 $4d$ BPS 态的 BPS 简并度 $\Omega(\gamma)$ 与增强 BPS 简并度 $\omega(\gamma,\overline{a})$。

这些技术使得可以计算 S 类中更高秩理论的 $4d$ BPS 谱。例子包括 22.2.4 节末尾提到的「五边形」Argyres–Douglas 理论、强耦合纯 $SU(3)$ SYM，以及九个自由超多重态的理论，如 [154, 第 8 节] 所讨论。许多其他重要例子随后在后续工作中得到处理，包括一般 Ar-

% [原文 p.395]
gyres–Douglas 理论 [156, 263]、$D$ 型与 $E$ 型\footnote{本书前页所画的图描绘的是纯 $SO(8)$ SYM 理论在其库仑分支原点处的谱网络。}的 S 类理论 [249]，以及 Minahan–Nemeschansky SCFT [206, 187]。

例如，$E_7$ Minahan–Nemeschansky 理论的 BPS 谱在 [187] 中通过用小质量参数 $m$ 形变该理论而得到研究：取亚纯三次微分
\[
\phi_2=\frac{mz\,\dd z^2}{(z^3-1)^2}\,,\qquad \phi_3=-\frac{\dd z^3}{(z^3-1)^2}
\]
来参数化库仑分支。用前几节概述的方法可以定义并研究形变后理论的谱网络 $\mathcal{W}_\vartheta$，并通过取极限 $m\to0$ 导出无质量理论的谱网络 $\mathcal{W}_\vartheta$。例如，可以证明当 $n$ 很大时渐近地有
\[
|\Omega(n\gamma_1)|\sim c\,n^{-\frac{5}{2}}\big(17+12\sqrt{2}\big)^n,
\]
其中 $c$ 是常数，并由此推出：与 $E_6$ Minahan–Nemeschansky 理论类似，$E_7$ 理论中携带「基本电荷的 $n$ 倍」这种电磁荷的 BPS 态，其指标总是 $(-1)^{n+1}n$ 的正整数倍。

% [Tutor Note: 渐近式中的底数可化简：$17+12\sqrt2=(3+2\sqrt2)^2$（因 $(3+2\sqrt2)^2=9+8+12\sqrt2$），而 $3+2\sqrt2=(1+\sqrt2)^2$，故 $17+12\sqrt2=(1+\sqrt2)^4$。因此 $|\Omega(n\gamma_1)|\sim c\,n^{-5/2}(1+\sqrt2)^{4n}$，其指数增长率 $\log(17+12\sqrt2)=4\log(1+\sqrt2)\approx3.525$ 与谱网络计数中常见的 $e^{\pi R|Z|}$ 型增长相容（$\log(1+\sqrt2)$ 是 silver ratio 的对数，即 Pell 数增长率）。核对：$n^{-5/2}$ 的幂次与「单圈精确 WKB 求和中 $n$ 个孤子链的贡献 $\sim n^{-3/2}$ 再乘以 $n^{-1}$（从积分度量得到）」的计数惯例相符，具体幂次依文献约定可能相差整数幂。]

\begin{remark}[22.33]
本节所述的谱网络从属于 Seiberg–Witten 全纯分歧覆叠 $\Sigma\to C_{\mathrm{UV}}$，因而有时称为 WKB 谱网络。然而，谱网络的概念可以推广到纯拓扑的设定。此时它们从属于不带解析结构的分歧覆叠 $\Sigma\to C$，称为拓扑谱网络或一般谱网络。
\end{remark}

尽管拓扑谱网络与 $4d$ $\mathcal{N}=2$ 理论中的 BPS 计数问题没有直接联系，它们仍保持其路径提升性质。具体而言，生成函数 $F(\wp,\vartheta)$ 的扭化同伦不变性使得一个称为非阿贝尔化的过程成为可能。给定 $\Sigma$ 上的阿贝尔平坦联络与 $C$ 上的谱网络，可以通过定义平行移动而在 $C$ 上构造非阿贝尔平坦联络。这一构造在 [154, 第 9 与第 10 节] 中针对 $C$ 上扭曲的 $\GL_n(\C)$-局部系统发展起来。其逆过程称为阿贝尔化。

最终，谱网络给出光滑黎曼曲面上扭曲、带装饰的非阿贝尔平坦联络模空间上的坐标卡。这里讨论的谱网络与路径提升规则适用于扭曲的带装饰局部系统，其模空间由 Fock–Goncharov $\mathcal{A}$-坐标描述。相反，第 III 部分特别是第 15 章所关注的略作修改的概念，则适用于带标架局部系统与 $X$-坐标。这一「带标架」版本的（非）阿贝尔化为带标架非阿贝尔平坦联络的模空间给出坐标卡，详见第 III 部分。

% [原文 p.396]
\begin{remark}[22.34]
研究 $4d$ $\mathcal{N}=2$ 理论中 BPS 谱的另一个基本方法是 BPS 箭图理论，它最早在 [15, 16] 中引入与研究（综述另见 [73]）。这一方法适用于具有箭图性质的理论，即在电荷格 $\Gamma$ 中存在一组生成元 $e_i$，使得对每个电荷为 $\gamma\in\Gamma$ 的 BPS 粒子，$\gamma$ 或 $-\gamma$ 落在由 $e_i$ 生成的正锥之内。

这类理论的 BPS 谱可以通过把表示论方法应用于与该理论相伴的 BPS 箭图来计算。BPS 箭图理论中的一个关键概念是变异，它是在 5.5 节簇代数簇的语境中引入的。特别地，所谓的变异方法为确定 BPS 谱提供了一个有效的框架。

人们曾尝试在谱网络理论与 BPS 箭图理论之间建立联系。例如，这类关系的一些探索可参见 [145, 143]。
\end{remark}
